\documentclass[12pt,a4paper]{article}
\usepackage[utf8]{inputenc}
\usepackage[T1]{fontenc}
\usepackage[spanish]{babel}
\usepackage{amsmath,amssymb,amsfonts}
\usepackage{graphicx}
\usepackage{xcolor}
\usepackage{hyperref}
\usepackage{geometry}
\usepackage{fancyhdr}
\usepackage{titlesec}
\usepackage{float}
\usepackage{caption}
\usepackage{booktabs}
\usepackage{longtable}
\usepackage{array}
\usepackage{verbatim}
\usepackage[most]{tcolorbox}
\geometry{margin=2.5cm}
\setlength{\headheight}{13.6pt}
% ============================================================
% PALETA DE COLORES 3Dosim
% ============================================================
\definecolor{cPrimary}{RGB}{30,64,124}      % azul oscuro titulos
\definecolor{cSecondary}{RGB}{70,130,180}   % azul claro subtitulos
\definecolor{cAccent}{RGB}{220,50,50}        % rojo AI Supervisor / alertas
\definecolor{cBg}{RGB}{245,245,245}          % gris muy claro fondo
\definecolor{cDarkBg}{RGB}{230,230,230}      % gris fondo tablas

% ============================================================
% FORMATO DE SECCIONES CON COLORES
% ============================================================
\titleformat*{\section}{\normalfont\Large\bfseries\color{cPrimary}}
\titleformat*{\subsection}{\normalfont\large\bfseries\color{cSecondary}}
\titleformat*{\subsubsection}{\normalfont\normalsize\bfseries\color{cPrimary}}
\pagestyle{fancy}
\fancyhf{}
\fancyhead[L]{\small 3Dosim --- Modulo 2}
\fancyhead[R]{\small Generacion de Input MCNP}
\fancyfoot[C]{\thepage}
\renewcommand{\headrulewidth}{0.4pt}
\newcommand{\vardef}[2]{\noindent\textbf{#1}: #2\\}
\setlength{\parskip}{0.5em}
\title{\textbf{Modulo 2: Generacion de Input MCNP}\\[0.5cm]
       \large Pipeline de Simulacion Monte Carlo para Dosimetria Hepatica\\[0.3cm]
       \normalsize Documentacion Tecnica --- Revision v4.0}
\author{Proyecto 3Dosim --- Dosimetria Hepatica con $^{90}$Y}
\date{Julio 2026}
% ============================================================
% CAJA COLOREADA PARA AI SUPERVISOR
% ============================================================
\tcbset{
  aiSupervisor/.style = {
    colback=cBg,
    colframe=cAccent,
    boxrule=0.7pt,
    arc=3pt,
    left=4mm,
    right=4mm,
    top=3mm,
    bottom=3mm,
    fonttitle=\bfseries\large,
    title=Control de Calidad (AI Supervisor),
    coltitle=white,
    colbacktitle=cAccent,
    attach boxed title to top left={yshift=-3mm,xshift=5mm},
    boxed title style={colback=cAccent,arc=2pt}
  }
}
\begin{document}
\maketitle
\thispagestyle{empty}
\newpage


% ============================================================
% RESUMEN DEL MODULO
% ============================================================
\section*{Resumen}
\addcontentsline{toc}{section}{Resumen}

Este modulo forma parte del pipeline completo \textbf{3Dosim} para dosimetria hepatica con Yttrium-90.
Su objetivo es transformar imagenes DICOM (CT anatomico + PET funcional) en una \textbf{labelmap dosimetrica}
-- un volumen 3D donde cada voxel se clasifica segun su tipo de tejido siguiendo la nomenclatura ICRP-110/ICRU-44.
Esta labelmap constituye la entrada fundamental para el Modulo 2, donde se genera el input de simulacion Monte Carlo (MCNP).

\vspace{0.5em}
\noindent
\textbf{Pregunta que resuelve:} \textit{Cual es la forma mas precisa y automatica de identificar y clasificar cada tejido del paciente para poder calcular la dosis absorbida por voxel?}

\vspace{0.3em}
\noindent
\textbf{Resultado principal:} Labelmap NIfTI/NRRD con 53+ materiales ICRP-110 indexados, lista para MCNP.


% ============================================================
% GLOSARIO DE ACRONIMOS
% ============================================================
\section*{Glosario de Acronimos}
\addcontentsline{toc}{section}{Glosario de Acronimos}

\begin{description}
    \item[ACE] A Compact ENDF --- formato de datos nucleares.
    \item[Bq] Becquerel --- unidad de actividad radiactiva (1 desintegracion/s).
    \item[CDF] Cumulative Distribution Function.
    \item[CORA/CORB/CORC] Coordinate Axis bins --- ejes de la grilla TMESH.
    \item[CT] Computed Tomography.
    \item[CTME] Cutoff Time --- tiempo maximo de simulacion.
    \item[DE4/DF4] Dose Energy / Dose Function --- conversion fluencia a dosis.
    \item[DICOM] Digital Imaging and Communications in Medicine.
    \item[*F8] Pulse Height Tally --- detector tipo pulso.
    \item[F6] Energy Deposition Tally.
    \item[FILL] Tarjeta MCNP para rellenar celda con lattice.
    \item[GBq] Gigabecquerel ($10^9$ Bq).
    \item[HU] Hounsfield Units --- unidades de atenuacion en TC.
    \item[ICRP] International Commission on Radiological Protection.
    \item[ICRU] International Commission on Radiation Units.
    \item[IJK] Indices de voxel (columna, fila, slice).
    \item[IMP] Importance (importancia de celda en MCNP).
    \item[LAT] Lattice (grilla regular de celdas en MCNP).
    \item[LIKE n BUT] Sintaxis MCNP para reutilizar celda con modificaciones.
    \item[LPS] Left-Posterior-Superior (sistema de coordenadas MCNP).
    \item[MCNP] Monte Carlo N-Particle (codigo de transporte, LANL).
    \item[MCTAL] Archivo de salida de MCNP con resultados de tallies.
    \item[MIRD] Medical Internal Radiation Dose.
    \item[MODE] Modalidad de transporte (P: foton, E: electron).
    \item[MRML] Medical Reality Markup Language (formato interno de Slicer).
    \item[NRRD] Nearly Raw Raster Data (formato de volumen).
    \item[NPS] Number of Particle Histories.
    \item[PDF] Probability Density Function.
    \item[PEDEP] Particle Energy DEPosition.
    \item[PET] Positron Emission Tomography.
    \item[RAS] Right-Anterior-Superior (sistema de coordenadas Slicer).
    \item[RLE] Run-Length Encoding (compresion de geometria).
    \item[RPP] Rectangular Parallelepiped (caja ortogonal en MCNP).
    \item[SDEF] Source Definition (tarjeta de fuente en MCNP).
    \item[SIRT] Selective Internal Radiation Therapy.
    \item[SUV] Standardized Uptake Value.
    \item[TACE] Transarterial Chemoembolization.
    \item[TMESH] Tally MESH (grilla de dosis en MCNP).
    \item[TS] TotalSegmentator (segmentacion anatomica).
    \item[U] Universe (numero de universo en lattice).
    \item[VOI] Volume of Interest.
    \item[WGT] Weight (peso de la particula en MCNP).
    \item[XS] Cross Section (seccion eficaz).
\end{description}

\newpage

% ============================================================
% TABLA DE CONTENIDOS
% ============================================================
\tableofcontents
\newpage

% ============================================================
% 1. VISION GENERAL
% ============================================================
\section{Vision General --- Arquitectura del Modulo 2}
\label{sec:vision}

\subsection*{Acronimos usados en esta seccion}
\begin{tabular}{ll}
\toprule
\textbf{Acronimo} & \textbf{Significado} \\
\midrule
MCNP & Monte Carlo N-Particle (codigo de transporte, LANL) \\
TMESH & Tally MESH (grilla de dosis en MCNP) \\
SDEF & Source Definition (tarjeta de fuente en MCNP) \\
RLE & Run-Length Encoding (compresion de geometria) \\
.mrb & Medical Reality Bundle (escena 3D Slicer) \\
VOI & Volume of Interest \\
HU & Hounsfield Units \\
ICRP & International Commission on Radiological Protection \\
ICRU & International Commission on Radiation Units \\
RPP & Rectangular Parallelepiped \\
TS & TotalSegmentator \\
\bottomrule
\end{tabular}

\subsection{Contexto Clinico}

La dosimetria personalizada en terapias con radionuclidos ($^{90}$Y, $^{177}$Lu, $^{131}$I) requiere transportar particulas a traves de una geometria voxelizada derivada de la TC del paciente. El Modulo 2 toma la escena segmentada generada por el Modulo 1 (labelmap + transforms + volumenes) y produce un archivo de entrada listo para MCNP (.i) que describe:

\begin{itemize}
    \item La geometria voxelizada del paciente (lattice + RLE).
    \item La composicion y densidad de cada material (6 tejidos base + 53 tejidos ICRP-110).
    \item La fuente radiactiva (isotopo, actividad, distribucion espacial).
    \item Los detectores (tallies) para estimar dosis absorbida.
\end{itemize}

El archivo generado se entrega al Modulo 3 para ejecutar MCNP y procesar resultados.

\subsection{Pipeline de 6 Pasos}

El pipeline se orquesta desde \texttt{pipeline\_mod2.py} y ejecuta 6 pasos en secuencia estricta:

\begin{verbatim}
   +---------------------------------------------------------------+
   |                   pipeline_mod2.py                             |
   |                                                                |
   |   check_slicer()  ->  load_scene()  ->  scan_nodes()           |
   |                                                                |
   |   validate_prereqs()  ->  generate_mcnp()  ->  validate()      |
   +---------------------------------------------------------------+
\end{verbatim}

\begin{enumerate}
    \item \textbf{check\_slicer()}: Verifica que 3D Slicer este disponible con el modulo SlicerDosim cargado.
    \item \textbf{load\_scene(mrb\_path)}: Carga el archivo .mrb que contiene la TC, el labelmap, las transformaciones espaciales (RAS a voxel) y los volumenes ROI.
    \item \textbf{scan\_nodes()}: Recorre el scene tree de 3D Slicer e identifica \texttt{volume\_node}, \texttt{labelmap\_node} y \texttt{voi\_nodes}.
    \item \textbf{validate\_prereqs()}: Verifica que no haya valores nulos o inconsistentes antes de pasar a generacion.
    \item \textbf{generate\_mcnp(params)}: Paso central que invoca \texttt{mcnp\_generator.py} con los parametros del tratamiento y construye las 12 secciones del archivo .i.
    \item \textbf{validate\_mcnp\_input(ruta\_i)}: Post-procesamiento que verifica tamano, keywords, coordenadas, tarjetas de corte y archivos auxiliares.
\end{enumerate}

\subsection{Parametros de Configuracion}

El pipeline se configura mediante \texttt{pipeline\_config.jsonc}:

\begin{verbatim}
{
  "scene_output_dir": "C:/MAT/3Dosim/ai-pipe/imagenes",
  "mcnp": {
    "isotope": "Y90",
    "n_particles": 1000000,
    "refine_hu": false,
    "flip_rows": false,
    "flip_z": false
  },
  "tumor": {
    "mode": "ts_liver_lesions",
    "synthetic_radius_mm": 10.0,
    "create_healthy_liver": true
  }
}
\end{verbatim}

\subsection{Relacion con Modulo 1 y Modulo 3}

\begin{verbatim}
  Modulo 1                Modulo 2                Modulo 3
  ---------               ---------               ---------
  Escena .mrb      --->   pipeline_mod2.py  --->  MCNP6
  (TC + labelmap           |                        |
   + VOIs)                 |-- check_slicer()       |
                           |-- load_scene()         |
                           |-- scan_nodes()         |
                           |-- validate_prereqs()   |
                           |-- generate_mcnp()      |
                           |     +---> paciente.i   |
                           +-- validate()     --->  mcnp6
                                                      |
                                                 Resultados
                                                 (mcnp_output)
\end{verbatim}

\subsection{Equivalencia MATLAB a Python}

\begin{center}
\begin{tabular}{ll}
\toprule
\textbf{MATLAB (v3.14)} & \textbf{Python (v4)} \\
\midrule
\texttt{modulo\_2/Generar\_Input\_MCNP} & \texttt{mcnp\_generator.py} (orquestador) \\
\quad crear geometria & \texttt{mcnp\_geometry.py} (RPP, lattice, RLE) \\
\quad asignar materiales & \texttt{mcnp\_materials.py} (TS\_SEGMENT\_MAP, ICRP-110) \\
\quad definir fuente & \texttt{mcnp\_source.py} (SDEF, espectro, .src) \\
\quad configurar tallies & \texttt{mcnp\_tallies.py} (TMESH, *F8, F6) \\
\quad escribir archivo .i & \texttt{generate()} $\to$ \texttt{paciente.i} \\
\bottomrule
\end{tabular}
\end{center}


\begin{tcolorbox}[aiSupervisor]

\begin{tcolorbox}[aiSupervisor]
\subsection*{Control de Calidad (AI Supervisor)}

\begin{table}[H]
\centering
\begin{tabular}{llc}
\toprule
\textbf{Verificacion} & \textbf{Metrica} & \textbf{Umbral} \\
\midrule
Archivo .i generado & Archivo no existe o $<$ 10 KB & Critico \\
Keywords MCNP presentes & Falta C, cells, surfaces, data & Critico \\
Coordenadas RPP vs TMESH & Rango TMESH fuera del RPP & Critico \\
Tarjetas CUT presentes & Modo P o E sin CUT:P/E & Critico \\
Materiales asignados & Voxel con indice 0 (sin material) & Critico \\
Archivo .src generado & No existe o esta vacio & Critico \\
\bottomrule
\end{tabular}
\end{table}

\newpage

% ============================================================
% 2. CARGA DE ESCENA
% ============================================================

\end{tcolorbox}


\begin{table}[H]
\centering
\begin{tabular}{llc}
\toprule
\textbf{Verificacion} & \textbf{Metrica} & \textbf{Umbral} \\
\midrule
Archivo .i generado & Archivo no existe o $<$ 10 KB & Critico \\
Keywords MCNP presentes & Falta C, cells, surfaces, data & Critico \\
Coordenadas RPP vs TMESH & Rango TMESH fuera del RPP & Critico \\
Tarjetas CUT presentes & Modo P o E sin CUT:P/E & Critico \\
Materiales asignados & Voxel con indice 0 (sin material) & Critico \\
Archivo .src generado & No existe o esta vacio & Critico \\
\bottomrule
\end{tabular}
\end{table}

\newpage

% ============================================================
% 2. CARGA DE ESCENA
% ============================================================

\end{tcolorbox}


\begin{table}[H]
\centering
\begin{tabular}{llc}
\toprule
\textbf{Verificacion} & \textbf{Metrica} & \textbf{Umbral} \\
\midrule
Archivo .i generado & Archivo no existe o $<$ 10 KB & Critico \\
Keywords MCNP presentes & Falta C, cells, surfaces, data & Critico \\
Coordenadas RPP vs TMESH & Rango TMESH fuera del RPP & Critico \\
Tarjetas CUT presentes & Modo P o E sin CUT:P/E & Critico \\
Materiales asignados & Voxel con indice 0 (sin material) & Critico \\
Archivo .src generado & No existe o esta vacio & Critico \\
\bottomrule
\end{tabular}
\end{table}

\newpage

% ============================================================
% 2. CARGA DE ESCENA
% ============================================================

\end{tcolorbox}


\begin{table}[H]
\centering
\begin{tabular}{llc}
\toprule
\textbf{Verificacion} & \textbf{Metrica} & \textbf{Umbral} \\
\midrule
Archivo .i generado & Archivo no existe o $<$ 10 KB & Critico \\
Keywords MCNP presentes & Falta C, cells, surfaces, data & Critico \\
Coordenadas RPP vs TMESH & Rango TMESH fuera del RPP & Critico \\
Tarjetas CUT presentes & Modo P o E sin CUT:P/E & Critico \\
Materiales asignados & Voxel con indice 0 (sin material) & Critico \\
Archivo .src generado & No existe o esta vacio & Critico \\
\bottomrule
\end{tabular}
\end{table}

\newpage

% ============================================================
% 2. CARGA DE ESCENA
% ============================================================
\section{Carga de Escena 3D Slicer}
\label{sec:carga}

\subsection*{Acronimos usados en esta seccion}
\begin{tabular}{ll}
\toprule
\textbf{Acronimo} & \textbf{Significado} \\
\midrule
.mrb & Medical Reality Bundle (escena comprimida de 3D Slicer) \\
RAS & Right-Anterior-Superior (sistema de coordenadas Slicer) \\
IJK & Indices de voxel (columna, fila, slice) \\
LPS & Left-Posterior-Superior (sistema MCNP) \\
VOI & Volume of Interest \\
MRML & Medical Reality Markup Language \\
NRRD & Nearly Raw Raster Data \\
\bottomrule
\end{tabular}

\subsection{Contexto}

El Modulo 1 produce una escena de 3D Slicer guardada como archivo .mrb. Este archivo contiene: el volumen de TC original en HU, el labelmap con segmentacion anatomica (TotalSegmentator + refinamiento medico), los volumenes ROI (contornos de higado, tumor, rinones) y las transformaciones espaciales que vinculan coordenadas RAS del mundo real con coordenadas IJK de voxel.

Cargar correctamente esta escena es el paso inicial del Modulo 2 y condiciona toda la generacion posterior.

\subsection{Archivo .mrb}

El formato .mrb es un ZIP que contiene:

\begin{verbatim}
escena.mrb/
+-- scene.mrml          (Arbol de nodos: volumenes, modelos, transforms)
+-- volumen_tc.nrrd     (TC en formato NRRD: espaciado, origen, datos)
+-- labelmap.nrrd       (Segmentacion: indices de material por voxel)
+-- higado.vtk          (Malla del contorno hepatico)
+-- tumor.vtk           (Malla del contorno tumoral)
+-- ...                 (Otros VOIs)
\end{verbatim}

\subsection{Carga desde pipeline\_mod2.py}

La funcion \texttt{load\_scene(mrb\_path)} usa la API de 3D Slicer:

\begin{verbatim}
def load_scene(mrb_path):
    slicer.util.loadScene(mrb_path)
    return {"status": "loaded", "path": mrb_path}
\end{verbatim}

\subsection{scan\_nodes() --- Identificar Nodos}

Una vez cargada la escena, \texttt{scan\_nodes()} recorre el arbol MRML:

\begin{center}
\begin{tabular}{lll}
\toprule
\textbf{Nodo} & \textbf{Tipo MRML} & \textbf{Proposito} \\
\midrule
\texttt{volume\_node} & \texttt{vtkMRMLScalarVolumeNode} & TC original (referencia espacial) \\
\texttt{labelmap\_node} & \texttt{vtkMRMLLabelMapVolumeNode} & Segmentacion con materiales \\
\texttt{voi\_nodes[]} & \texttt{vtkMRMLModelNode[]} & Mallas de VOIs \\
\bottomrule
\end{tabular}
\end{center}

Desde el \texttt{volume\_node} se extraen:

\begin{verbatim}
origen_ras = volume_node.GetOrigin()        # (x, y, z) en mm
espaciado  = volume_node.GetSpacing()       # (dx, dy, dz) en mm
dims       = volume_node.GetImageData().GetDimensions()  # (nx, ny, nz)
\end{verbatim}

\subsection{Coordenadas IJK a RAS}

La matriz de transformacion IJK$\to$RAS se obtiene con:

\begin{verbatim}
ijk_to_ras = vtk.vtkMatrix4x4()
volume_node.GetIJKToRASMatrix(ijk_to_ras)
\end{verbatim}

Esta matriz de $4 \times 4$ permite convertir cualquier indice de voxel $(i,j,k)$ a coordenadas RAS en mm:

$$
\begin{bmatrix} x_{RAS} \\ y_{RAS} \\ z_{RAS} \\ 1 \end{bmatrix}
= \mathbf{M}_{IJK \to RAS} \cdot \begin{bmatrix} i \\ j \\ k \\ 1 \end{bmatrix}
$$

\vardef{$\mathbf{M}_{IJK \to RAS}$}{Matriz de transformacion $4\times4$ de coordenadas voxel (IJK) a coordenadas mundo (RAS)}

\begin{center}
\begin{tabular}{llll}
\toprule
\textbf{Sistema} & \textbf{Eje 1} & \textbf{Eje 2} & \textbf{Eje 3} \\
\midrule
Slicer (RAS) & R (Right) $+$izq & A (Anterior) $+$atras & S (Superior) $+$arriba \\
MCNP (LPS) & L--R (Left$\to$Right) & P--A (Post$\to$Ant) & S--I (Sup$\to$Inf) \\
Voxel (IJK) & i (columna) & j (fila) & k (slice) \\
\bottomrule
\end{tabular}
\end{center}

La conversion de RAS a LPS requiere:

$$
x_{LPS} = -x_{RAS},\qquad y_{LPS} = -y_{RAS},\qquad z_{LPS} = z_{RAS}
$$

\vardef{$x_{LPS}$}{Coordenada X en sistema LPS (MCNP), igual a $-x_{RAS}$}
\vardef{$y_{LPS}$}{Coordenada Y en sistema LPS (MCNP), igual a $-y_{RAS}$}
\vardef{$z_{LPS}$}{Coordenada Z en sistema LPS (MCNP), igual a $z_{RAS}$}

Los flags \texttt{flip\_rows} y \texttt{flip\_z} del generador MCNP manejan la reorientacion de ejes entre ambos sistemas.

\subsection{Matriz IJK$\to$RAS en Detalle}

$$
\mathbf{M}_{IJK \to RAS} = \begin{bmatrix}
S_x R_{00} & S_y R_{01} & S_z R_{02} & O_x \\
S_x R_{10} & S_y R_{11} & S_z R_{12} & O_y \\
S_x R_{20} & S_y R_{21} & S_z R_{22} & O_z \\
0 & 0 & 0 & 1
\end{bmatrix}
$$

\vardef{$S_x, S_y, S_z$}{Espaciado del voxel en mm en cada direccion}
\vardef{$R$}{Matriz de rotacion $3\times3$ del volumen}
\vardef{$O_x, O_y, O_z$}{Origen en RAS del voxel $(0,0,0)$ en IJK}

\subsection{Ejemplo de Extraccion de Parametros}

Para un volumen tipico de TC abdominal:

\begin{verbatim}
origin_ras = (-120.5, -150.3, -245.0)   # mm
spacing    = (0.98, 0.98, 2.50)          # mm/voxel
dims       = (512, 512, 180)             # voxels

# Calcular RPP en LPS (para MCNP)
x_min = -(-120.5) - 0.98 * 512 / 2 = 120.5 - 250.88 = -130.38
x_max = -(-120.5) + 0.98 * 512 / 2 = 120.5 + 250.88 = 371.38
y_min = -(-150.3) - 0.98 * 512 / 2 = 150.3 - 250.88 = -100.58
y_max = -(-150.3) + 0.98 * 512 / 2 = 150.3 + 250.88 = 401.18
z_min = -245.0 - 2.50 * 180 / 2 = -245.0 - 225.0 = -470.0
z_max = -245.0 + 2.50 * 180 / 2 = -245.0 + 225.0 = -20.0
\end{verbatim}

\subsection{Validacion de Prerequisitos}

\texttt{validate\_prereqs()} ejecuta 4 verificaciones:

\begin{enumerate}
    \item \textbf{Volumen de TC presente}: \texttt{assert volume\_node is not None}
    \item \textbf{Labelmap cargado}: \texttt{assert labelmap\_node is not None}
    \item \textbf{Dimensiones compatibles}: \texttt{assert tc\_dims == lm\_dims}
    \item \textbf{VOIs requeridos}: \texttt{assert "Liver" in voi\_names}, etc.
\end{enumerate}

\subsection{Manejo de Errores}

\begin{center}
\begin{tabular}{lll}
\toprule
\textbf{Error} & \textbf{Causa} & \textbf{Accion} \\
\midrule
\texttt{FileNotFoundError} & Ruta .mrb invalida & Verificar selector de archivos \\
\texttt{ValueError} & Nodos incompatibles & Revisar version de Slicer \\
\texttt{AssertionError} & Labelmap sin materiales & Ejecutar Modulo 1 primero \\
\texttt{RuntimeError} & Slicer no disponible & Iniciar SlicerDosim \\
\bottomrule
\end{tabular}
\end{center}


\begin{tcolorbox}[aiSupervisor]

\begin{tcolorbox}[aiSupervisor]
\subsection*{Control de Calidad (AI Supervisor)}

\begin{table}[H]
\centering
\begin{tabular}{llc}
\toprule
\textbf{Verificacion} & \textbf{Metrica} & \textbf{Umbral} \\
\midrule
Escena cargada exitosamente & Error en \texttt{loadScene()} & Critico \\
TC y labelmap mismas dims & \texttt{tc\_dims $\neq$ lm\_dims} & Critico \\
VOIs requeridos presentes & Falta Liver/Tumor/Kidney & Critico \\
Matriz IJK$\to$RAS invertible & Determinante cero & Critico \\
Espaciado valido & $dx \leq 0$ o $dy \leq 0$ o $dz \leq 0$ & Critico \\
\bottomrule
\end{tabular}
\end{table}

\newpage

% ============================================================
% 3. ASIGNACION DE MATERIALES
% ============================================================

\end{tcolorbox}


\begin{table}[H]
\centering
\begin{tabular}{llc}
\toprule
\textbf{Verificacion} & \textbf{Metrica} & \textbf{Umbral} \\
\midrule
Escena cargada exitosamente & Error en \texttt{loadScene()} & Critico \\
TC y labelmap mismas dims & \texttt{tc\_dims $\neq$ lm\_dims} & Critico \\
VOIs requeridos presentes & Falta Liver/Tumor/Kidney & Critico \\
Matriz IJK$\to$RAS invertible & Determinante cero & Critico \\
Espaciado valido & $dx \leq 0$ o $dy \leq 0$ o $dz \leq 0$ & Critico \\
\bottomrule
\end{tabular}
\end{table}

\newpage

% ============================================================
% 3. ASIGNACION DE MATERIALES
% ============================================================

\end{tcolorbox}


\begin{table}[H]
\centering
\begin{tabular}{llc}
\toprule
\textbf{Verificacion} & \textbf{Metrica} & \textbf{Umbral} \\
\midrule
Escena cargada exitosamente & Error en \texttt{loadScene()} & Critico \\
TC y labelmap mismas dims & \texttt{tc\_dims $\neq$ lm\_dims} & Critico \\
VOIs requeridos presentes & Falta Liver/Tumor/Kidney & Critico \\
Matriz IJK$\to$RAS invertible & Determinante cero & Critico \\
Espaciado valido & $dx \leq 0$ o $dy \leq 0$ o $dz \leq 0$ & Critico \\
\bottomrule
\end{tabular}
\end{table}

\newpage

% ============================================================
% 3. ASIGNACION DE MATERIALES
% ============================================================

\end{tcolorbox}


\begin{table}[H]
\centering
\begin{tabular}{llc}
\toprule
\textbf{Verificacion} & \textbf{Metrica} & \textbf{Umbral} \\
\midrule
Escena cargada exitosamente & Error en \texttt{loadScene()} & Critico \\
TC y labelmap mismas dims & \texttt{tc\_dims $\neq$ lm\_dims} & Critico \\
VOIs requeridos presentes & Falta Liver/Tumor/Kidney & Critico \\
Matriz IJK$\to$RAS invertible & Determinante cero & Critico \\
Espaciado valido & $dx \leq 0$ o $dy \leq 0$ o $dz \leq 0$ & Critico \\
\bottomrule
\end{tabular}
\end{table}

\newpage

% ============================================================
% 3. ASIGNACION DE MATERIALES
% ============================================================
\section{Asignacion de Materiales}
\label{sec:materiales}

\subsection*{Acronimos usados en esta seccion}
\begin{tabular}{ll}
\toprule
\textbf{Acronimo} & \textbf{Significado} \\
\midrule
HU & Hounsfield Units (unidades de atenuacion radiologica) \\
ICRP & International Commission on Radiological Protection \\
ICRU & International Commission on Radiation Units \\
TS & TotalSegmentator \\
TOI & Tissue of Interest \\
Z & Numero atomico \\
RLE & Run-Length Encoding \\
UD & Universe/Detail (tarjeta UD en MCNP lattice) \\
\bottomrule
\end{tabular}

\subsection{Contexto}

Cada voxel en el labelmap contiene un indice numerico que representa el tejido biologico en esa posicion. Para que MCNP pueda transportar particulas correctamente, cada indice debe mapear a:

\begin{enumerate}
    \item Una composicion elemental (fraccion en peso de cada isotopo).
    \item Una densidad (en g/cm$^3$).
    \item Un rango de HU de referencia (opcional, para validacion).
\end{enumerate}

El sistema asigna materiales usando dos fuentes de datos en paralelo: 6 tejidos base definidos localmente en \texttt{mcnp\_materials.py} y 53 materiales del ICRP-110 cargados desde \texttt{tissue\_config.json}.

\subsection{Indices de Tejidos Base (Phantom, indices 1--100)}

\begin{center}
\begin{tabular}{llccc}
\toprule
\textbf{Indice} & \textbf{Tejido} & $\rho$ (g/cm$^3$) & \textbf{Rango HU} & \textbf{ID MCNP} \\
\midrule
1 & Aire (Dry near sea level) & 0.001205 & $-$1024 a $-$900 & 1 \\
30 & Tejido blando (ICRU 44) & 1.06 & $-$200 a 150 & 2 \\
50 & Pulmon (ICRP 110) & 0.50 & $-$900 a $-$200 & 4 \\
80 & Hueso (ICRP 110) & 1.60 & 150 a 2000 & 5 \\
90 & Higado (ICRP 110) & 1.05 & 50 a 150 & 3 \\
100 & Tumor & 1.05 & 50 a 150 & 6 \\
\bottomrule
\end{tabular}
\end{center}

\subsection{Indices ICRP-110 (Completos, indices 101--153)}

El archivo \texttt{tissue\_config.json} contiene 53 materiales adicionales del ICRP-110. Cada entrada tiene:

\begin{verbatim}
{
  "index": 110,
  "name": "Hueso_vertebral",
  "name_en": "Trabecular_bone",
  "hu_range": [150, 300],
  "density_gcm3": 1.35,
  "mcnp_material": {
    "id": 10,
    "composition": {
      "1000": 0.074,
      "6000": 0.254,
      ...
    }
  }
}
\end{verbatim}

La composicion se expresa como fraccion en peso de cada isotopo (Zzzzz AAAA, por ejemplo $1000 = {}^1$H, $6000 = {}^{12}$C).

\subsection{Carga desde tissue\_config.json}

\texttt{TissueConfig.load\_config()} en \texttt{config.py} implementa el patron singleton:

\begin{verbatim}
class TissueConfig:
    _instance = None
    tissues: list[TissueEntry]
    index_map: dict[int, TissueEntry]
    material_map: dict[int, TissueEntry]

    @classmethod
    def load_config(cls, path=None):
        if cls._instance is None:
            cls._instance = cls()
            with open(path or "Resources/Config/tissue_config.json") as f:
                data = json.load(f)
            for t in data["tissues"]:
                cls._instance.index_map[t["index"]] = t
                cls._instance.material_map[t["mcnp_material"]["id"]] = t
        return cls._instance
\end{verbatim}

\subsection{Segment Map (TotalSegmentator $\to$ Indices Locales)}

La constante \texttt{TS\_SEGMENT\_MAP} en \texttt{mcnp\_materials.py} mapea los nombres de segmentos de TotalSegmentator a los indices del phantom:

\begin{verbatim}
TS_SEGMENT_MAP = {
    "skin": 30,                    # Tejido_blando
    "subcutaneous_fat": 30,
    "bone": 80,                    # Hueso
    "rib": 80,
    "lung_upper_lobe_left": 50,    # Pulmon
    "lung_lower_lobe_left": 50,
    "lung_upper_lobe_right": 50,
    "lung_middle_lobe_right": 50,
    "lung_lower_lobe_right": 50,
    "liver": 90,                   # Higado
    # ...
}
\end{verbatim}

Los nombres que no aparecen en el mapa se asignan por defecto a Tejido\_blando (30). El tumor (100) se asigna manualmente.

\subsection{Refinamiento por HU}

Cuando \texttt{refine\_hu=True}, el generador analiza la TC original en cada voxel para ajustar el material segun el valor de HU:

\begin{verbatim}
                  +------------------+
                  |  Voxel (i,j)     |
                  |  HU = valor      |
                  +--------+---------+
                           |
                  +--------+---------+
                  |                   v
            HU en rango          HU fuera de rango
            del segmento         del segmento
                  |                   |
                  v                   v
            conserva indice     recalcular por HU
                                       |
                           +-----------+-----------+
                           |                       v
                     HU ~ Aire               HU ~ Tejido
                     indice=1                indice=30
\end{verbatim}

Categorias de umbral:

\begin{center}
\begin{tabular}{llc}
\toprule
\textbf{Condicion} & \textbf{Nuevo material} & \textbf{Rango HU} \\
\midrule
$HU < -900$ & Aire (1) & $-$1024 a $-$900 \\
$HU > 150$ & Hueso (80) & 150 a 2000 \\
$-900 \leq HU \leq 150$ & Tejido blando (30) & $-$900 a 150 \\
\bottomrule
\end{tabular}
\end{center}

\subsection{Materiales en el Input MCNP}

Cada material unico genera una tarjeta Mn en el bloque de datos:

\begin{verbatim}
M1    6000 0.000124  7000 0.755268  8000 0.231481  18000 0.012827  $ Aire
M2    1000 0.105     6000 0.143     7000 0.034     8000 0.708     $ Tejido_blando
      11000 0.002    15000 0.003    16000 0.003    17000 0.002
      19000 0.003
M3    1000 0.102     6000 0.131     7000 0.031     8000 0.724     $ Higado
...
\end{verbatim}

Las densidades se aplican en la celda correspondiente:

\begin{verbatim}
100 1 -0.001205 -1 2 3 ...  $ Celda de aire (material 1, rho=0.001205)
200 2 -1.06      -4 5 6 ...  $ Celda de tejido blando (material 2, rho=1.06)
\end{verbatim}


\begin{tcolorbox}[aiSupervisor]

\begin{tcolorbox}[aiSupervisor]
\subsection*{Control de Calidad (AI Supervisor)}

\begin{table}[H]
\centering
\begin{tabular}{llc}
\toprule
\textbf{Verificacion} & \textbf{Metrica} & \textbf{Umbral} \\
\midrule
Mapa TS\_SEGMENT\_MAP completo & Segmento sin mapeo & Critico \\
Indices labelmap validos (1 a 153) & Indice 0 o $>$ 153 & Critico \\
Densidad fisica & $\rho \leq 0$ o $\rho \geq 20$ g/cm$^3$ & Critico \\
Composicion suma $\sim 1.0$ & $|\Sigma \text{fracciones} - 1| > 10^{-4}$ & Critico \\
Refinamiento HU plausible & $> 50\%$ voxeles recategorizados & Advertencia \\
Archivo tissue\_config.json cargable & Parse error & Critico \\
\bottomrule
\end{tabular}
\end{table}

\newpage

% ============================================================
% 4. GEOMETRIA VOXELIZADA MCNP
% ============================================================

\end{tcolorbox}


\begin{table}[H]
\centering
\begin{tabular}{llc}
\toprule
\textbf{Verificacion} & \textbf{Metrica} & \textbf{Umbral} \\
\midrule
Mapa TS\_SEGMENT\_MAP completo & Segmento sin mapeo & Critico \\
Indices labelmap validos (1 a 153) & Indice 0 o $>$ 153 & Critico \\
Densidad fisica & $\rho \leq 0$ o $\rho \geq 20$ g/cm$^3$ & Critico \\
Composicion suma $\sim 1.0$ & $|\Sigma \text{fracciones} - 1| > 10^{-4}$ & Critico \\
Refinamiento HU plausible & $> 50\%$ voxeles recategorizados & Advertencia \\
Archivo tissue\_config.json cargable & Parse error & Critico \\
\bottomrule
\end{tabular}
\end{table}

\newpage

% ============================================================
% 4. GEOMETRIA VOXELIZADA MCNP
% ============================================================

\end{tcolorbox}


\begin{table}[H]
\centering
\begin{tabular}{llc}
\toprule
\textbf{Verificacion} & \textbf{Metrica} & \textbf{Umbral} \\
\midrule
Mapa TS\_SEGMENT\_MAP completo & Segmento sin mapeo & Critico \\
Indices labelmap validos (1 a 153) & Indice 0 o $>$ 153 & Critico \\
Densidad fisica & $\rho \leq 0$ o $\rho \geq 20$ g/cm$^3$ & Critico \\
Composicion suma $\sim 1.0$ & $|\Sigma \text{fracciones} - 1| > 10^{-4}$ & Critico \\
Refinamiento HU plausible & $> 50\%$ voxeles recategorizados & Advertencia \\
Archivo tissue\_config.json cargable & Parse error & Critico \\
\bottomrule
\end{tabular}
\end{table}

\newpage

% ============================================================
% 4. GEOMETRIA VOXELIZADA MCNP
% ============================================================

\end{tcolorbox}


\begin{table}[H]
\centering
\begin{tabular}{llc}
\toprule
\textbf{Verificacion} & \textbf{Metrica} & \textbf{Umbral} \\
\midrule
Mapa TS\_SEGMENT\_MAP completo & Segmento sin mapeo & Critico \\
Indices labelmap validos (1 a 153) & Indice 0 o $>$ 153 & Critico \\
Densidad fisica & $\rho \leq 0$ o $\rho \geq 20$ g/cm$^3$ & Critico \\
Composicion suma $\sim 1.0$ & $|\Sigma \text{fracciones} - 1| > 10^{-4}$ & Critico \\
Refinamiento HU plausible & $> 50\%$ voxeles recategorizados & Advertencia \\
Archivo tissue\_config.json cargable & Parse error & Critico \\
\bottomrule
\end{tabular}
\end{table}

\newpage

% ============================================================
% 4. GEOMETRIA VOXELIZADA MCNP
% ============================================================
\section{Geometria Voxelizada MCNP}
\label{sec:geometria}

\subsection*{Acronimos usados en esta seccion}
\begin{tabular}{ll}
\toprule
\textbf{Acronimo} & \textbf{Significado} \\
\midrule
RPP & Rectangular Parallelepiped (caja ortogonal en MCNP) \\
LIKE n BUT & Sintaxis MCNP para reutilizar celda con modificaciones \\
U & Universe (numero de universo en lattice) \\
FILL & Tarjeta MCNP para rellenar celda con lattice \\
LAT & Lattice (grilla regular de celdas en MCNP) \\
RLE & Run-Length Encoding (codificacion por longitud de carrera) \\
TRCL & Transformacion de coordenadas en celda MCNP \\
\bottomrule
\end{tabular}

\subsection{Estrategia de Geometria Voxelizada}

MCNP no soporta voxeles nativamente, pero si lattices. La geometria se estructura en 3 niveles:

\begin{verbatim}
                     +---------------------------+
                     |   Celda mundo (999)        |
                     |   RPP contenedor           |
                     |   FILL = universo 1        |
                     +-------------+--------------+
                                   |
                     +-------------+--------------+
                     |   Celda lattice (1)        |
                     |   LAT=1                    |
                     |   RPP = volumen             |
                     |   FILL = i:nx j:ny k:nz    |
                     +-------------+--------------+
                                   |
                  +----------------+----------------+
                  v                v                v
            Celda voxel       Celda voxel       Celda voxel
            (100 n)           (100 n)           (100 n)
            LIKE 100 BUT      LIKE 100 BUT      LIKE 100 BUT
            MAT = material    MAT = material    MAT = material
\end{verbatim}

\subsection{Celda Mundo (Envolvente)}

Se calculan las coordenadas extremas del volumen en LPS:

$$
\begin{aligned}
x_{min} &= -origin_{RAS,x} - \frac{espaciado_x \cdot n_x}{2} \\
x_{max} &= -origin_{RAS,x} + \frac{espaciado_x \cdot n_x}{2} \\
y_{min} &= -origin_{RAS,y} - \frac{espaciado_y \cdot n_y}{2} \\
y_{max} &= -origin_{RAS,y} + \frac{espaciado_y \cdot n_y}{2} \\
z_{min} &= origin_{RAS,z} - \frac{espaciado_z \cdot n_z}{2} \\
z_{max} &= origin_{RAS,z} + \frac{espaciado_z \cdot n_z}{2}
\end{aligned}
$$

\vardef{$x_{min}, x_{max}$}{Limites del RPP en X (LPS) en cm}
\vardef{$origin_{RAS,x}$}{Origen del volumen en X segun Slicer (RAS) en mm}
\vardef{$espaciado_x$}{Espaciado del voxel en X (mm/voxel)}
\vardef{$n_x$}{Numero de voxeles en direccion X}

\begin{verbatim}
999 0 -1000 $ Celda mundo (vacio exterior)
1000 RPP x_min x_max y_min y_max z_min z_max
\end{verbatim}

\subsection{Celda Lattice}

Se usa el operador \texttt{LAT=1} (grilla ortogonal):

\begin{verbatim}
1 0 -1000 U=1 LAT=1
    FILL=-nx:2*nx -ny:2*ny -nz:2*nz
    1:n_rev(ny*nz), 1:n_rev((ny-1)*nz), ..., 1:3, 1:2, 1:1
\end{verbatim}

El rango del FILL se define con:

$$
i \in [0, n_x),\qquad j \in [0, n_y),\qquad k \in [0, n_z)
$$

\vardef{$i, j, k$}{Indices de voxel en el lattice MCNP}
\vardef{$n_x, n_y, n_z$}{Dimensiones del volumen en voxeles}

\subsection{Celdas Voxel con LIKE n BUT}

\texttt{mcnp\_geometry.py} genera celdas voxel usando \texttt{LIKE n BUT}:

\begin{verbatim}
100 1 -0.001205 U=1 $ Celda base: aire
     $ sin RPP -- la celda se define por el lattice
101 LIKE 100 BUT MAT=2 $ Tejido blando
    RHO=1.06
102 LIKE 100 BUT MAT=3 $ Higado
    RHO=1.05
...
\end{verbatim}

Cada celda voxel unica se define una sola vez. La matriz FILL del lattice asigna que celda (material) esta en cada posicion $(i,j,k)$.

\subsection{Compresion RLE}

Para reducir el tamano del archivo de entrada, el generador implementa Run-Length Encoding (RLE) sobre el plano k:

\begin{verbatim}
Sin RLE:
    FILL=-1:1 -1:1 0:0
        100 101 102
        101 100 101
        102 101 100

Con RLE:
    En lugar de repetir 100, 100, 100, 100, 100 se escribe:
    5R 100
\end{verbatim}

El algoritmo recorre cada fila del plano k y reemplaza secuencias del mismo material por el factor de repeticion:

\begin{verbatim}
def comprimir_rle(filas, n_cols):
    resultado = []
    for fila in filas:
        count = 0
        prev = None
        linea = []
        for val in fila:
            if val == prev:
                count += 1
            else:
                if prev is not None:
                    if count == 1:
                        linea.append(str(prev))
                    else:
                        linea.append(f"{count}R {prev}")
                prev = val
                count = 1
        if count == 1:
            linea.append(str(prev))
        else:
            linea.append(f"{count}R {prev}")
        resultado.append(" ".join(linea))
    return "\n".join(resultado)
\end{verbatim}

\subsection{Reorientacion de Ejes (flip\_rows y flip\_z)}

Los flags \texttt{flip\_rows} y \texttt{flip\_z} controlan la orientacion del volumen.

\textbf{flip\_rows}: Invierte el orden de las filas (eje j) en el lattice:

$$
j_{MCNP} = (n_y - 1) - j_{RAS}
$$

\vardef{$j_{MCNP}$}{Indice de fila en coordenadas MCNP}
\vardef{$j_{RAS}$}{Indice de fila en coordenadas RAS (Slicer)}

\textbf{flip\_z}: Invierte el orden de los slices (eje k):

\begin{verbatim}
if flip_z:
    slices = slices[::-1]  # slice 0 pasa a ser el ultimo
\end{verbatim}

La transformacion completa es:

$$
\begin{aligned}
i_{MCNP} &= i_{RAS} \\
j_{MCNP} &= \begin{cases}
j_{RAS} & \text{si } \neg flip\_rows \\
(n_y - 1) - j_{RAS} & \text{si } flip\_rows
\end{cases} \\
k_{MCNP} &= \begin{cases}
k_{RAS} & \text{si } \neg flip\_z \\
(n_z - 1) - k_{RAS} & \text{si } flip\_z
\end{cases}
\end{aligned}
$$

\vardef{$i_{MCNP}, j_{MCNP}, k_{MCNP}$}{Indices de voxel en el sistema MCNP}
\vardef{$flip\_rows$}{Flag booleano que indica si se invierte el eje de filas}
\vardef{$flip\_z$}{Flag booleano que indica si se invierte el eje de slices}

\subsection{Precision de Cuantizacion}

Las coordenadas de las superficies RPP se cuantizan a 1 micra de precision:

$$
sx_q = \frac{\text{round}(sx_{mm} \cdot 1000)}{1000}
$$

\vardef{$sx_q$}{Espaciado cuantizado a 1 micra (mm)}
\vardef{$sx_{mm}$}{Espaciado original en mm}

Las dimensiones totales del volumen resultan:

$$
x_m = \frac{n_x \cdot sx_q}{10}\text{ cm},\qquad
y_m = \frac{n_y \cdot sy_q}{10}\text{ cm},\qquad
z_m = \frac{n_z \cdot sz_q}{10}\text{ cm}
$$

\vardef{$x_m$}{Dimension total en X del RPP en cm}
\vardef{$sx_q$}{Espaciado cuantizado en X (mm)}
\vardef{$n_x$}{Numero de voxeles en X}

\subsection{Superficies Especiales}

Ademas del RPP contenedor, se define una esfera para debugging:

\begin{verbatim}
1 RPP x_min x_max y_min y_max z_min z_max   $ Caja contenedora
2 RPP -sx_q/2 sx_q/2 -sy_q/2 sy_q/2 -sz_q/2 sz_q/2  $ Voxel unitario
650 SO 15     $ Esfera de 15 cm de radio (debug)
\end{verbatim}

\subsection{Estructura Completa en el .i}

\begin{verbatim}
999 0 -1000        imp:p=1 imp:e=1
1000 RPP -150 150 -200 200 -300 0

1 0 -1000 u=1 lat=1
    fill=-5:2*5 -4:2*4 -3:2*3
    ...

100 1 -0.001205 u=1  $ Aire
101 2 -1.06     u=1  $ Tejido_blando
102 3 -1.05     u=1  $ Higado
103 4 -0.50     u=1  $ Pulmon
104 5 -1.60     u=1  $ Hueso
105 6 -1.05     u=1  $ Tumor
...
153 54 -DENSIDAD u=1  $ ICRP-110 ultimo material
\end{verbatim}

\subsection{Consideraciones de Rendimiento}

\begin{center}
\begin{tabular}{ll}
\toprule
\textbf{Factor} & \textbf{Impacto} \\
\midrule
Tamano del .i & Lineal con $n_{voxels}$ unicos (RLE reduce $\sim 10\times$) \\
Tiempo de carga MCNP & Proporcional al tamano del FILL \\
Memoria MCNP & 1--2 GB para volumenes tipicos ($512^3$) \\
Velocidad de transporte & Independiente de la geometria (depende de materiales) \\
\bottomrule
\end{tabular}
\end{center}


\begin{tcolorbox}[aiSupervisor]

\begin{tcolorbox}[aiSupervisor]
\subsection*{Control de Calidad (AI Supervisor)}

\begin{table}[H]
\centering
\begin{tabular}{llc}
\toprule
\textbf{Verificacion} & \textbf{Metrica} & \textbf{Umbral} \\
\midrule
RPP del lattice contiene todo el body & Voxeles fuera del RPP & Critico \\
FILL tiene $nx \times ny \times nz$ elementos & Cuenta incorrecta & Critico \\
Todos los materiales tienen celda U=1 & Falta celda voxel & Critico \\
RLE mantiene integridad de materiales & Suma repeticiones $\neq nx \times ny$ & Critico \\
flip\_rows/flip\_z consistentes con image\_origin & Error de registro espacial & Critico \\
\bottomrule
\end{tabular}
\end{table}

\newpage

% ============================================================
% 5. DEFINICION DE LA FUENTE
% ============================================================

\end{tcolorbox}


\begin{table}[H]
\centering
\begin{tabular}{llc}
\toprule
\textbf{Verificacion} & \textbf{Metrica} & \textbf{Umbral} \\
\midrule
RPP del lattice contiene todo el body & Voxeles fuera del RPP & Critico \\
FILL tiene $nx \times ny \times nz$ elementos & Cuenta incorrecta & Critico \\
Todos los materiales tienen celda U=1 & Falta celda voxel & Critico \\
RLE mantiene integridad de materiales & Suma repeticiones $\neq nx \times ny$ & Critico \\
flip\_rows/flip\_z consistentes con image\_origin & Error de registro espacial & Critico \\
\bottomrule
\end{tabular}
\end{table}

\newpage

% ============================================================
% 5. DEFINICION DE LA FUENTE
% ============================================================

\end{tcolorbox}


\begin{table}[H]
\centering
\begin{tabular}{llc}
\toprule
\textbf{Verificacion} & \textbf{Metrica} & \textbf{Umbral} \\
\midrule
RPP del lattice contiene todo el body & Voxeles fuera del RPP & Critico \\
FILL tiene $nx \times ny \times nz$ elementos & Cuenta incorrecta & Critico \\
Todos los materiales tienen celda U=1 & Falta celda voxel & Critico \\
RLE mantiene integridad de materiales & Suma repeticiones $\neq nx \times ny$ & Critico \\
flip\_rows/flip\_z consistentes con image\_origin & Error de registro espacial & Critico \\
\bottomrule
\end{tabular}
\end{table}

\newpage

% ============================================================
% 5. DEFINICION DE LA FUENTE
% ============================================================

\end{tcolorbox}


\begin{table}[H]
\centering
\begin{tabular}{llc}
\toprule
\textbf{Verificacion} & \textbf{Metrica} & \textbf{Umbral} \\
\midrule
RPP del lattice contiene todo el body & Voxeles fuera del RPP & Critico \\
FILL tiene $nx \times ny \times nz$ elementos & Cuenta incorrecta & Critico \\
Todos los materiales tienen celda U=1 & Falta celda voxel & Critico \\
RLE mantiene integridad de materiales & Suma repeticiones $\neq nx \times ny$ & Critico \\
flip\_rows/flip\_z consistentes con image\_origin & Error de registro espacial & Critico \\
\bottomrule
\end{tabular}
\end{table}

\newpage

% ============================================================
% 5. DEFINICION DE LA FUENTE
% ============================================================
\section{Definicion de la Fuente (Source)}
\label{sec:fuente}

\subsection*{Acronimos usados en esta seccion}
\begin{tabular}{ll}
\toprule
\textbf{Acronimo} & \textbf{Significado} \\
\midrule
SDEF & Source Definition (tarjeta MCNP para definir la fuente) \\
WGT & Weight (peso de la particula) \\
ERG & Energy (energia de emision) \\
POS & Position (posicion de la fuente) \\
.src & Archivo externo con posiciones de fuente \\
XS & Cross Section (seccion eficaz) \\
PDF & Probability Density Function \\
CDF & Cumulative Distribution Function \\
ACE & A Compact ENDF (formato de datos nucleares) \\
\bottomrule
\end{tabular}

\subsection{Contexto}

En dosimetria de radionuclidos, la fuente de radiacion esta distribuida en el volumen del paciente (no es un haz externo). El $^{90}$Y, $^{177}$Lu o $^{131}$I se administra al paciente y se acumula selectivamente en tejidos. La fuente MCNP debe representar esta distribucion volumetrica con la energia espectral correcta.

\subsection{Espectro de Energias --- Y-90 ($^{90}$Y)}

El $^{90}$Y decae por emision beta pura a $^{90}$Zr (estable). La energia maxima es:

$$
E_{\beta,\text{max}}^{^{90}\text{Y}} = 2.28\ \text{MeV}
$$

\vardef{$E_{\beta,\text{max}}$}{Energia maxima de la emision beta del $^{90}$Y (2.28 MeV)}

La energia media es:

$$
\bar{E}_\beta = \frac{1}{3} E_{\beta,\text{max}} \approx 0.935\ \text{MeV}
$$

\vardef{$\bar{E}_\beta$}{Energia media del espectro beta del $^{90}$Y}

El espectro se modela con 2 grupos de energia:

\begin{center}
\begin{tabular}{lcc}
\toprule
\textbf{Grupo} & \textbf{Rango (MeV)} & \textbf{Probabilidad} \\
\midrule
1 & 0.0 -- 0.5 & 0.30 \\
2 & 0.5 -- 2.28 & 0.70 \\
\bottomrule
\end{tabular}
\end{center}

La funcion de densidad de probabilidad (PDF) del espectro beta teorico de Fermi:

$$
\frac{dN}{dE} = C \cdot F(Z, E) \cdot p \cdot E \cdot (E_{\max} - E)^2
$$

\vardef{$C$}{Constante de normalizacion del espectro beta}
\vardef{$F(Z, E)$}{Funcion de Fermi (correccion coulombiana)}
\vardef{$p$}{Momento del electron ($p = \sqrt{E^2 - m_e^2 c^4}$)}
\vardef{$E$}{Energia cinetica del electron}
\vardef{$E_{\max}$}{Energia maxima del decaimiento beta (2.28 MeV para $^{90}$Y)}
\vardef{$Z$}{Numero atomico del nucleo hijo ($^{90}$Zr, $Z=40$)}

\subsection{Espectro Lu-177 y I-131}

\begin{center}
\begin{tabular}{llcc}
\toprule
\textbf{Isotopo} & \textbf{Grupo} & \textbf{Rango (MeV)} & \textbf{Probabilidad} \\
\midrule
$^{177}$Lu & 1 & 0.0 -- 0.2 & 0.15 \\
$^{177}$Lu & 2 & 0.2 -- 0.5 & 0.60 \\
$^{177}$Lu & 3 & 0.5 -- 0.8 & 0.25 \\
\midrule
$^{131}$I ($e^-$) & 1 & 0.0 -- 0.6 & 0.90 \\
$^{131}$I ($\gamma$) & 2 & 0.364 & 0.10 \\
\bottomrule
\end{tabular}
\end{center}

\subsection{Tarjeta SDEF}

La tarjeta SDEF generada para fuente volumetrica con archivo .src externo:

\begin{verbatim}
SDEF PAR=D1 ERG=FPAR=D6 X=D2 Y=D3 Z=D4 CELL=D5 WGT=1.00033788800193
SI2 H 0. sx_cm
SP2 D 0 1
SI5 L (u lc[ix iy iz] lw)
SP5 peso=1/n_active
read file Y90cel3D.src
\end{verbatim}

\vardef{$WGT$}{Peso de la particula en MCNP (1.00033788800193 para normalizacion)}

Los parametros clave son:

\begin{description}
    \item[PAR=D1] Tipo de particula (1 = foton, 2 = electron).
    \item[ERG=FPAR=D6] Energia desde distribucion D6 (espectro).
    \item[X=D2 Y=D3 Z=D4] Posicion muestreada desde distribuciones uniformes.
    \item[CELL=D5] Celda de la fuente.
    \item[WGT=1.00033788800193] Peso constante para normalizacion.
\end{description}

Distribuciones uniformes en voxel:

\begin{verbatim}
SI2 H 0. sx_cm
SP2 D 0 1
\end{verbatim}

El espectro se define con:

\begin{verbatim}
SI1 L 0.9357 2.2807  $ energias discretas
SP1 1.0 0.0          $ probabilidades
\end{verbatim}

\subsection{Archivo de Fuente Externo (.src)}

Para fuentes distribuidas, el generador produce un archivo donde cada linea contiene las coordenadas de emision de una particula:

\begin{verbatim}
# Archivo de fuente 3Dosim
# Isotopo: Y90
# N particulas: 1000000
# Formato: x(mm) y(mm) z(mm) energia(MeV) peso
-45.2 30.1 -12.5 0.93 1.0
-44.8 29.7 -12.1 1.45 1.0
...
\end{verbatim}

\subsection{Muestreo desde Labelmap}

Cuando no se dispone de una malla VOI, el generador distribuye las particulas uniformemente sobre los voxeles tumorales:

\begin{verbatim}
def sample_positions_in_labelmap(labelmap, tumor_index=100, n_particles=1e6):
    mask = labelmap == tumor_index
    voxel_indices = argwhere(mask)
    n_voxels = len(voxel_indices)
    selected = voxel_indices[randint(0, n_voxels, n_particles)]
    positions = selected + uniform(-0.5, 0.5, selected.shape)
    return positions * espaciado + origen_lps
\end{verbatim}

\subsection{Actividad y Peso}

El peso de cada particula (WGT) representa su fraccion de la actividad total:

$$
WGT_i = \frac{A_{total}}{N_{particulas}}
$$

\vardef{$A_{total}$}{Actividad total administrada en Bq}
\vardef{$WGT_i$}{Peso de la particula $i$ en Bq/historia}

Para MCNP, el peso se normaliza para que la suma de todos los pesos sea 1.0:

$$
\dot{D} = \frac{A_{total}}{N_{particulas}} \cdot \sum_{i=1}^{N} D_i
$$

\vardef{$\dot{D}$}{Tasa de dosis absorbida}

\subsection{Tarjetas Relacionadas}

\begin{verbatim}
CUT:P J 1e-3    $ Corte de energia para fotones (1 keV)
CUT:E J 1e-3    $ Corte de energia para electrones (1 keV)
MODE P E        $ Transportar fotones y electrones
PHYS:E 1 0 0 0 1 0 0.99  $ Habilitar e- secundarios
\end{verbatim}

\subsection{Parametros de Fuente}

\begin{center}
\begin{tabular}{lll}
\toprule
\textbf{Parametro} & \textbf{Valor tipico} & \textbf{Dependencia} \\
\midrule
Isotopo & Y90 / Lu177 / I131 & Input del usuario \\
N particulas & $1\times10^6$ a $1\times10^7$ & Exactitud deseada \\
Actividad & 0.5--5 GBq (Y90 SIRT) & Informe medico \\
Energia media & 0.935 MeV (Y90) & Isotopo \\
Alcance medio $\beta$ & 2.5 mm (Y90 en tejido) & Isotopo + densidad \\
Volumen fuente & Tumor / Higado & Segmentacion \\
\bottomrule
\end{tabular}
\end{center}


\begin{tcolorbox}[aiSupervisor]

\begin{tcolorbox}[aiSupervisor]
\subsection*{Control de Calidad (AI Supervisor)}

\begin{table}[H]
\centering
\begin{tabular}{llc}
\toprule
\textbf{Verificacion} & \textbf{Metrica} & \textbf{Umbral} \\
\midrule
Espectro definido para isotopo & Falta entrada en SPECTRA dict & Critico \\
Archivo .src generado (si aplica) & No existe o 0 particulas & Critico \\
SDEF tiene ERG y WGT consistentes & Mismatch entre parametros & Critico \\
Peso suma $\approx 1.0$ & $|\Sigma WGT - 1| > 0.01$ & Critico \\
Particulas dentro del VOI & Coordenadas fuera de la malla & Critico \\
MODE coincide con la particula & Modo P sin electrones para beta & Critico \\
\bottomrule
\end{tabular}
\end{table}

\newpage

% ============================================================
% 6. TALLIES Y DETECTORES
% ============================================================

\end{tcolorbox}


\begin{table}[H]
\centering
\begin{tabular}{llc}
\toprule
\textbf{Verificacion} & \textbf{Metrica} & \textbf{Umbral} \\
\midrule
Espectro definido para isotopo & Falta entrada en SPECTRA dict & Critico \\
Archivo .src generado (si aplica) & No existe o 0 particulas & Critico \\
SDEF tiene ERG y WGT consistentes & Mismatch entre parametros & Critico \\
Peso suma $\approx 1.0$ & $|\Sigma WGT - 1| > 0.01$ & Critico \\
Particulas dentro del VOI & Coordenadas fuera de la malla & Critico \\
MODE coincide con la particula & Modo P sin electrones para beta & Critico \\
\bottomrule
\end{tabular}
\end{table}

\newpage

% ============================================================
% 6. TALLIES Y DETECTORES
% ============================================================

\end{tcolorbox}


\begin{table}[H]
\centering
\begin{tabular}{llc}
\toprule
\textbf{Verificacion} & \textbf{Metrica} & \textbf{Umbral} \\
\midrule
Espectro definido para isotopo & Falta entrada en SPECTRA dict & Critico \\
Archivo .src generado (si aplica) & No existe o 0 particulas & Critico \\
SDEF tiene ERG y WGT consistentes & Mismatch entre parametros & Critico \\
Peso suma $\approx 1.0$ & $|\Sigma WGT - 1| > 0.01$ & Critico \\
Particulas dentro del VOI & Coordenadas fuera de la malla & Critico \\
MODE coincide con la particula & Modo P sin electrones para beta & Critico \\
\bottomrule
\end{tabular}
\end{table}

\newpage

% ============================================================
% 6. TALLIES Y DETECTORES
% ============================================================

\end{tcolorbox}


\begin{table}[H]
\centering
\begin{tabular}{llc}
\toprule
\textbf{Verificacion} & \textbf{Metrica} & \textbf{Umbral} \\
\midrule
Espectro definido para isotopo & Falta entrada en SPECTRA dict & Critico \\
Archivo .src generado (si aplica) & No existe o 0 particulas & Critico \\
SDEF tiene ERG y WGT consistentes & Mismatch entre parametros & Critico \\
Peso suma $\approx 1.0$ & $|\Sigma WGT - 1| > 0.01$ & Critico \\
Particulas dentro del VOI & Coordenadas fuera de la malla & Critico \\
MODE coincide con la particula & Modo P sin electrones para beta & Critico \\
\bottomrule
\end{tabular}
\end{table}

\newpage

% ============================================================
% 6. TALLIES Y DETECTORES
% ============================================================
\section{Tallies y Detectores MCNP}
\label{sec:tallies}

\subsection*{Acronimos usados en esta seccion}
\begin{tabular}{ll}
\toprule
\textbf{Acronimo} & \textbf{Significado} \\
\midrule
TMESH & Tally MESH (grilla 3D de dosis en MCNP) \\
CORA/CORB/CORC & Coordinate Axis bins (ejes de la grilla TMESH) \\
DE4/DF4 & Dose Energy / Dose Function (conversion fluencia $\to$ dosis) \\
*F8 & Pulse Height Tally (detector tipo pulso) \\
F6 & Energy Deposition Tally \\
PEDEP & Particle Energy DEPosition \\
Gy & Gray (J/kg, unidad SI de dosis absorbida) \\
\bottomrule
\end{tabular}

\subsection{Contexto}

Las tallies (detectores) en MCNP son el mecanismo para estimar la dosis absorbida en el paciente. 3Dosim utiliza tres tipos complementarios:

\begin{itemize}
    \item \textbf{TMESH rmesh1:e pedep}: grilla 3D de dosis para mapas volumetricos.
    \item \textbf{*F8}: detectores puntuales en organos de interes (higado, tumor, rinones).
    \item \textbf{F6:e}: energia depositada total para verificacion de conservacion.
\end{itemize}

\subsection{TMESH --- Tally de Grilla 3D}

TMESH es un tally tipo 1 (fluencia) sobre una malla cartesiana. La definicion en el input MCNP:

\begin{verbatim}
TMESH
  RMESH1:PEDEP PEDEP
  CORA1  x_min  x_max  nx
  CORB1  y_min  y_max  ny
  CORC1  z_min  z_max  nz
  ENDSD
\end{verbatim}

\begin{center}
\begin{tabular}{lll}
\toprule
\textbf{Parametro} & \textbf{Significado} & \textbf{Valor tipico} \\
\midrule
$x_{min}$, $x_{max}$ & Limites grilla en X (mm) & $-150$ a $150$ \\
$y_{min}$, $y_{max}$ & Limites grilla en Y (mm) & $-200$ a $200$ \\
$z_{min}$, $z_{max}$ & Limites grilla en Z (mm) & $-300$ a $0$ \\
$nx$, $ny$, $nz$ & Numero de intervalos por eje & 128 a 256 \\
\bottomrule
\end{tabular}
\end{center}

\subsubsection{BUG CORREGIDO --- TMESH nx vs nx-1}

Versiones anteriores del generador usaban:

\begin{verbatim}
CORA1 x_min x_max nx i  $ INCORRECTO
\end{verbatim}

La correccion cambio a:

\begin{verbatim}
CORA1 x_min x_max nx  $ CORRECTO -- MCNP espera nx intervalos
\end{verbatim}

\subsection{Conversion de Fluencia a Dosis (DE4/DF4)}

El tally RMESH1:PEDEP usa un par DE4/DF4 para convertir fluencia de electrones (MeV/cm$^2$) a dosis (MeV/g). Los 11 puntos de la conversion:

\begin{verbatim}
DE4  0.01 0.03 0.05 0.07 0.10 0.20 0.40 0.60 1.00 2.00 2.28
DF4  1.0  1.0  1.0  1.0  1.0  1.0  1.0  1.0  1.0  1.0  1.0
\end{verbatim}

\vardef{$DE4$}{Vector de energias para conversion fluencia $\to$ dosis (MeV)}
\vardef{$DF4$}{Vector de factores de conversion (MeV/g)/(MeV/cm$^2$)}

La conversion responde a la formula general:

$$
H(E) = \frac{\mu_{en}(E)}{\rho} \cdot \Phi(E) \cdot E
$$

\vardef{$H(E)$}{Dosis en energia $E$ (MeV/g)}
\vardef{$\mu_{en}(E)/\rho$}{Coeficiente masico de absorcion de energia (cm$^2$/g)}
\vardef{$\Phi(E)$}{Fluencia de particulas en energia $E$ (particulas/cm$^2$)}

Post-procesamiento: la dosis en MeV/g se convierte a Gy:

$$
D_{Gy} = D_{MeV/g} \cdot \frac{1.602176634 \times 10^{-13}\ \text{J/MeV}}{10^{-3}\ \text{kg/g}}
$$

$$
D_{Gy} = D_{MeV/g} \times 1.602176634 \times 10^{-10}
$$

\vardef{$D_{Gy}$}{Dosis absorbida en Gray (J/kg)}
\vardef{$D_{MeV/g}$}{Dosis en unidades MCNP (MeV/g)}

\subsection{*F8 --- Detectores Puntuales por Organo}

Los tallies *F8 (pulse height) proporcionan el espectro de energia depositada en una celda. 3Dosim define *F8 por cada VOI con muestreo aleatorio de voxeles:

\begin{verbatim}
*F8:P  (90 lc [x1 y1 z1])  $ Tally F8 en celda higado
*F18:P (100 lc [x2 y2 z2])  $ Tally F8 en celda tumor
\end{verbatim}

La notacion \texttt{lc} permite especificar una lista de celdas que conforman el detector.

La desviacion estandar relativa (RSD) para *F8:

$$
RSD = \frac{1}{\sqrt{N_{eventos}}}
$$

\vardef{$RSD$}{Desviacion estandar relativa del tally *F8}
\vardef{$N_{eventos}$}{Numero de eventos detectados en el tally}

Para alcanzar $RSD < 5\%$:

$$
N_{eventos} > \left(\frac{1}{0.05}\right)^2 = 400
$$

\subsection{F6:e --- Energia Depositada Total}

El tally F6 mide la energia depositada por unidad de masa (MeV/g):

\begin{verbatim}
F6:P  body_cell  $ Energia depositada en todo el paciente
\end{verbatim}

Se usa como verificacion de conservacion:

$$
E_{depositada}^{F6} \approx E_{emitida}^{SDEF}
$$

\subsection{Tarjetas de Corte (CUT)}

Controlan la energia minima de simulacion:

\begin{verbatim}
CUT:P J 1e-3    $ Corte para fotones (1 keV)
CUT:E J 1e-3    $ Corte para electrones (1 keV)
\end{verbatim}

El parametro \texttt{J} indica que es un corte en energia (no en tiempo).

\subsection{Conversion de Unidades}

\begin{center}
\begin{tabular}{lll}
\toprule
\textbf{Unidad MCNP} & \textbf{Unidad SI} & \textbf{Factor} \\
\midrule
MeV/g & Gy & $1.602176634 \times 10^{-10}$ \\
MeV/cm$^2$ & -- & (fluencia, se multiplica por $\mu_{en}/\rho$) \\
shakes & s & $10^{-8}$ \\
\bottomrule
\end{tabular}
\end{center}

\subsection{Resumen de Tallies}

\begin{center}
\begin{tabular}{llll}
\toprule
\textbf{Tally} & \textbf{ID} & \textbf{Tipo} & \textbf{Proposito} \\
\midrule
RMESH1:PEDEP & 1 & Grilla 3D & Mapa de dosis volumetrico \\
DE4/DF4 & 4 & Conversion & Fluencia $\to$ dosis (11 puntos) \\
*F8:P & 8 & Puntual & Dosis en organo especifico \\
*F18:P & 18 & Puntual & Dosis en tumor \\
F6:P & 6 & Volumetrica & Energia depositada total \\
\bottomrule
\end{tabular}
\end{center}

\subsection{Parametros Globales de Simulacion}

\begin{verbatim}
NPS  10000000     $ 10 millones de historias
RAND STRIDE=1111152917 GEN=2 SEED=random  $ Semilla reproducible
\end{verbatim}

\vardef{$NPS$}{Numero de historias de particulas en la simulacion MCNP}
\vardef{$RAND$}{Tarjeta de control de numeros aleatorios}

\subsection{Verificacion de Conservacion de Energia}

$$
E_{emitida} = N_{particulas} \times \bar{E}_{\beta}
$$

\vardef{$E_{emitida}$}{Energia total emitida por la fuente (MeV)}

$$
E_{depositada} = \sum_{voxeles} D_{voxel} \times \rho_{voxel} \times V_{voxel}
$$

\vardef{$E_{depositada}$}{Energia total depositada en el volumen (MeV)}
\vardef{$\rho_{voxel}$}{Densidad del voxel (g/cm$^3$)}
\vardef{$V_{voxel}$}{Volumen de un voxel ($dx \cdot dy \cdot dz$)}

$$
\text{Ratio} = \frac{E_{depositada}}{E_{emitida}} \approx 1.0
$$

El ratio debe estar entre 0.95 y 1.05.


\begin{tcolorbox}[aiSupervisor]

\begin{tcolorbox}[aiSupervisor]
\subsection*{Control de Calidad (AI Supervisor)}

\begin{table}[H]
\centering
\begin{tabular}{llc}
\toprule
\textbf{Verificacion} & \textbf{Metrica} & \textbf{Umbral} \\
\midrule
TMESH cubre todo el volumen RPP & Rango TMESH $<$ RPP & Critico \\
CORA/CORB/CORC definidos & Falta alguno de los 3 ejes & Critico \\
nx, ny, nz enteros $>$ 0 & $nx \leq 0$ o no entero & Critico \\
DE4/DF4 tienen 11 puntos cada uno & Distinta cantidad de puntos & Critico \\
*F8 asocia celda existente & Celda del detector no definida & Critico \\
CUT:P y CUT:E presentes & Modo P o E sin CUT & Critico \\
ENDSD cierra TMESH & Falta ENDSD & Critico \\
\bottomrule
\end{tabular}
\end{table}

\newpage

% ============================================================
% 7. GENERACION DEL ARCHIVO .i
% ============================================================

\end{tcolorbox}


\begin{table}[H]
\centering
\begin{tabular}{llc}
\toprule
\textbf{Verificacion} & \textbf{Metrica} & \textbf{Umbral} \\
\midrule
TMESH cubre todo el volumen RPP & Rango TMESH $<$ RPP & Critico \\
CORA/CORB/CORC definidos & Falta alguno de los 3 ejes & Critico \\
nx, ny, nz enteros $>$ 0 & $nx \leq 0$ o no entero & Critico \\
DE4/DF4 tienen 11 puntos cada uno & Distinta cantidad de puntos & Critico \\
*F8 asocia celda existente & Celda del detector no definida & Critico \\
CUT:P y CUT:E presentes & Modo P o E sin CUT & Critico \\
ENDSD cierra TMESH & Falta ENDSD & Critico \\
\bottomrule
\end{tabular}
\end{table}

\newpage

% ============================================================
% 7. GENERACION DEL ARCHIVO .i
% ============================================================

\end{tcolorbox}


\begin{table}[H]
\centering
\begin{tabular}{llc}
\toprule
\textbf{Verificacion} & \textbf{Metrica} & \textbf{Umbral} \\
\midrule
TMESH cubre todo el volumen RPP & Rango TMESH $<$ RPP & Critico \\
CORA/CORB/CORC definidos & Falta alguno de los 3 ejes & Critico \\
nx, ny, nz enteros $>$ 0 & $nx \leq 0$ o no entero & Critico \\
DE4/DF4 tienen 11 puntos cada uno & Distinta cantidad de puntos & Critico \\
*F8 asocia celda existente & Celda del detector no definida & Critico \\
CUT:P y CUT:E presentes & Modo P o E sin CUT & Critico \\
ENDSD cierra TMESH & Falta ENDSD & Critico \\
\bottomrule
\end{tabular}
\end{table}

\newpage

% ============================================================
% 7. GENERACION DEL ARCHIVO .i
% ============================================================

\end{tcolorbox}


\begin{table}[H]
\centering
\begin{tabular}{llc}
\toprule
\textbf{Verificacion} & \textbf{Metrica} & \textbf{Umbral} \\
\midrule
TMESH cubre todo el volumen RPP & Rango TMESH $<$ RPP & Critico \\
CORA/CORB/CORC definidos & Falta alguno de los 3 ejes & Critico \\
nx, ny, nz enteros $>$ 0 & $nx \leq 0$ o no entero & Critico \\
DE4/DF4 tienen 11 puntos cada uno & Distinta cantidad de puntos & Critico \\
*F8 asocia celda existente & Celda del detector no definida & Critico \\
CUT:P y CUT:E presentes & Modo P o E sin CUT & Critico \\
ENDSD cierra TMESH & Falta ENDSD & Critico \\
\bottomrule
\end{tabular}
\end{table}

\newpage

% ============================================================
% 7. GENERACION DEL ARCHIVO .i
% ============================================================
\section{Generacion del Archivo .i}
\label{sec:generacion}

\subsection*{Acronimos usados en esta seccion}
\begin{tabular}{ll}
\toprule
\textbf{Acronimo} & \textbf{Significado} \\
\midrule
NPS & Number of Particles (numero de historias) \\
MODE & Modalidad de transporte (P: foton, E: electron) \\
IMP & Importance (importancia de celda) \\
PRDMP & Print Dump (control de impresion) \\
VOID & Tarjeta VOID (espacio vacio sin transporte) \\
RFL & Right, Left, Free (formato MCNP) \\
\bottomrule
\end{tabular}

\subsection{Contexto}

La generacion del input MCNP es el paso central del Modulo 2. Recibe los datos de la escena procesada (labelmap, TC, VOIs, parametros del tratamiento) y produce un archivo de texto con formato MCNP6 (.i) listo para ejecutar.

\subsection{Flujo de Generacion}

\begin{verbatim}
  Parametros del tratamiento
  (isotopo, n_particles, refine_hu, flip_rows, flip_z)
         |
         v
  mcnp_generator.generate()
         |
         +-- 1. Cabecera (comentarios + metadata)
         +-- 2. Bloque de celdas (cell builder)
         +-- 3. Bloque de superficies (surface builder)
         +-- 4. Tarjetas de datos (data builder)
         +-- 5. Materiales (material builder)
         +-- 6. Fuente SDEF (source builder)
         +-- 7. Tallies TMESH (tally builder)
         +-- 8. Tallies *F8 (detector builder)
         +-- 9. Tally F6
         +-- 10. Print
         +-- 11. Archivo .src externo
         +-- 12. Corte de geometria (geocut)
\end{verbatim}

\subsection{Parametros de Entrada}

\begin{verbatim}
params = {
    "isotope": "Y90",        # Y90, Lu177, I131
    "n_particles": 1000000,  # 1x10^6 historias
    "refine_hu": False,      # Refinar materiales por HU
    "flip_rows": False,      # Invertir filas RAS->MCNP
    "flip_z": False,         # Invertir slices
    "image_origin": None,    # Origen LPS opcional
}
\end{verbatim}

\subsection{Las 12 Secciones en Detalle}

\subsubsection{Seccion 1: Cabecera}

\begin{verbatim}
Mensaje de bloqueo (1 linea)

C ================================================================
C Archivo de entrada MCNP generado por 3Dosim v4.0
C ================================================================
C Paciente   : {patient_id}
C Isotopo    : {isotope}
C Particulas : {n_particles}
C Fecha      : {timestamp}
C Refine HU  : {refine_hu}
C Flip rows  : {flip_rows}
C Flip Z     : {flip_z}
C ================================================================
\end{verbatim}

\subsubsection{Seccion 2: Bloque de Celdas}

Define las celdas voxel usando \texttt{LIKE n BUT} y \texttt{LAT=1}:

\begin{verbatim}
C ===== CELDAS =====
C Celda mundo
999 0 -1000  IMP:P=1 IMP:E=1  $ Exterior (vacio)
C
C Celda lattice (voxels)
1 0 -1000 U=1 LAT=1
   FILL=-{nx}:{2*nx} -{ny}:{2*ny} -{nz}:{2*nz}
   1 1 1 1 2 2 2 3 3 3 3 4 ...
C
C Celdas voxel (materiales)
{v} LIKE 1 BUT MAT={mid} RHO=-{rho} U={v}
\end{verbatim}

\subsubsection{Seccion 3: Bloque de Superficies}

Define el RPP contenedor:

\begin{verbatim}
C ===== SUPERFICIES =====
1000 RPP {x_min} {x_max} {y_min} {y_max} {z_min} {z_max}
\end{verbatim}

\subsubsection{Seccion 4: Tarjetas de Datos}

\begin{verbatim}
C ===== DATOS =====
MODE P E
IMP:P 1  $ Importancia foton = 1 en todas las celdas
IMP:E 1  $ Importancia electron = 1 en todas las celdas
CUT:P J 1e-3  $ Corte de energia para fotones
CUT:E J 1e-3  $ Corte de energia para electrones
\end{verbatim}

\subsubsection{Seccion 5: Materiales}

\begin{verbatim}
C ===== MATERIALES =====
M1   6000 0.000124  7000 0.755268  8000 0.231481  18000 0.012827  $ Aire
M2   1000 0.105  6000 0.143  7000 0.034  8000 0.708  $ Tejido_blando
     11000 0.002  15000 0.003  16000 0.003  17000 0.002
     19000 0.003
M3   1000 0.102  6000 0.131  7000 0.031  8000 0.724  $ Higado
...
M6   1000 0.102  6000 0.131  7000 0.031  8000 0.724  $ Tumor
\end{verbatim}

\subsubsection{Seccion 6: Fuente SDEF}

\begin{verbatim}
C ===== FUENTE =====
SDEF PAR=D1 ERG=FPAR=D6 X=D2 Y=D3 Z=D4 CELL=D5
      WGT=1.00033788800193
SI2 H 0. sx_cm
SP2 D 0 1
SI5 L (u lc[ix iy iz] lw)
SP5 peso=1/n_active
read file Y90cel3D.src
\end{verbatim}

\subsubsection{Seccion 7: Tallies TMESH}

\begin{verbatim}
C ===== TMESH =====
TMESH
  RMESH1:PEDEP PEDEP
  CORA1 {x_min} {x_max} {nx}
  CORB1 {y_min} {y_max} {ny}
  CORC1 {z_min} {z_max} {nz}
  DE4 0.01 0.03 0.05 0.07 0.10 0.20 0.40 0.60 1.00 2.00 2.28
  DF4 1.0 1.0 1.0 1.0 1.0 1.0 1.0 1.0 1.0 1.0 1.0
  ENDSD
\end{verbatim}

\subsubsection{Seccion 8: Tallies *F8}

\begin{verbatim}
C ===== *F8 =====
*F8:P  (90 lc [x1 y1 z1])  $ Higado
*F18:P (100 lc [x2 y2 z2])  $ Tumor
\end{verbatim}

\subsubsection{Seccion 9: Tally F6}

\begin{verbatim}
C ===== F6 =====
F6:P {body_cell}  $ Energia depositada en todo el paciente
\end{verbatim}

\subsubsection{Seccion 10: Print}

\begin{verbatim}
C ===== PRINT =====
PRINT
PRDMP 2J 1
\end{verbatim}

\subsubsection{Seccion 11: Archivo de Fuente Externo}

Se genera un archivo \texttt{.src} con las posiciones muestreadas:

\begin{verbatim}
{output_dir}/{patient_id}.src
\end{verbatim}

\subsubsection{Seccion 12: Corte de Geometria (Opcional)}

Solo se incluye si la variable de entorno \texttt{DOSIM\_DEBUG\_GEOM} esta activa:

\begin{verbatim}
C ===== GEOCUT =====
VOID
\end{verbatim}

\subsection{Formato del Archivo}

El formato MCNP limita cada linea a 72 caracteres. El generador respeta esta restriccion:

\begin{verbatim}
{linea:72s}  $ Formato con padding a 72 caracteres
\end{verbatim}

\subsection{Validacion Post-Generacion}

Despues de escribir el archivo, \texttt{validate\_mcnp\_input()} ejecuta 5 verificaciones:

\textbf{Verificacion 1: Tamano minimo}

\begin{verbatim}
assert os.path.getsize(path) > 10_000, "Archivo demasiado pequeno"
\end{verbatim}

\textbf{Verificacion 2: Keywords obligatorias}

\begin{verbatim}
keywords = ["C ", "M", "SDEF", "TMESH", "CORA", "CORB", "CORC",
            "MODE", "CUT", "PRINT", "FILL", "LAT", "RPP"]
\end{verbatim}

\textbf{Verificacion 3: Coincidencia RPP vs TMESH}

Extraer limites y verificar que TMESH este dentro de RPP.

\textbf{Verificacion 4: Tarjetas CUT}

Si el modo incluye P, debe existir CUT:P; si incluye E, debe existir CUT:E.

\textbf{Verificacion 5: Archivo .src generado}

Si $n_{particulas} > 0$, el archivo .src debe existir y no estar vacio.

\subsection{Archivos de Salida}

\begin{center}
\begin{tabular}{lll}
\toprule
\textbf{Archivo} & \textbf{Contenido} & \textbf{Tamano tipico} \\
\midrule
\texttt{paciente.i} & Input MCNP completo & 5--100 MB \\
\texttt{paciente.src} & Posiciones de fuente & 10--100 MB \\
\texttt{paciente.dat} & Espectro de energias & 1--10 KB \\
\texttt{paciente.debug\_geom.txt} & Debug geometrico (si DEBUG) & 1--5 KB \\
\bottomrule
\end{tabular}
\end{center}


\begin{tcolorbox}[aiSupervisor]

\begin{tcolorbox}[aiSupervisor]
\subsection*{Control de Calidad (AI Supervisor)}

\begin{table}[H]
\centering
\begin{tabular}{llc}
\toprule
\textbf{Verificacion} & \textbf{Metrica} & \textbf{Umbral} \\
\midrule
Archivo .i generado $>$ 10 KB & Tamano $\leq$ 10 KB & Critico \\
Todas las keywords MCNP presentes & Keyword faltante & Critico \\
Coordenadas TMESH dentro del RPP & TMESH fuera de limites & Critico \\
Tarjetas CUT para modos activos & Modo sin CUT & Critico \\
Todas las celdas tienen material $>$ 0 & Celda con material 0 & Critico \\
FILL tiene $nx \times ny \times nz$ elementos & Cuenta incorrecta & Critico \\
Material M definido para cada celda & Material referenciado sin M & Critico \\
Archivo .src existe y tiene contenido & Archivo faltante o vacio & Critico \\
\bottomrule
\end{tabular}
\end{table}

\newpage

% ============================================================
% 8. PIPELINE COMPLETO
% ============================================================

\end{tcolorbox}


\begin{table}[H]
\centering
\begin{tabular}{llc}
\toprule
\textbf{Verificacion} & \textbf{Metrica} & \textbf{Umbral} \\
\midrule
Archivo .i generado $>$ 10 KB & Tamano $\leq$ 10 KB & Critico \\
Todas las keywords MCNP presentes & Keyword faltante & Critico \\
Coordenadas TMESH dentro del RPP & TMESH fuera de limites & Critico \\
Tarjetas CUT para modos activos & Modo sin CUT & Critico \\
Todas las celdas tienen material $>$ 0 & Celda con material 0 & Critico \\
FILL tiene $nx \times ny \times nz$ elementos & Cuenta incorrecta & Critico \\
Material M definido para cada celda & Material referenciado sin M & Critico \\
Archivo .src existe y tiene contenido & Archivo faltante o vacio & Critico \\
\bottomrule
\end{tabular}
\end{table}

\newpage

% ============================================================
% 8. PIPELINE COMPLETO
% ============================================================

\end{tcolorbox}


\begin{table}[H]
\centering
\begin{tabular}{llc}
\toprule
\textbf{Verificacion} & \textbf{Metrica} & \textbf{Umbral} \\
\midrule
Archivo .i generado $>$ 10 KB & Tamano $\leq$ 10 KB & Critico \\
Todas las keywords MCNP presentes & Keyword faltante & Critico \\
Coordenadas TMESH dentro del RPP & TMESH fuera de limites & Critico \\
Tarjetas CUT para modos activos & Modo sin CUT & Critico \\
Todas las celdas tienen material $>$ 0 & Celda con material 0 & Critico \\
FILL tiene $nx \times ny \times nz$ elementos & Cuenta incorrecta & Critico \\
Material M definido para cada celda & Material referenciado sin M & Critico \\
Archivo .src existe y tiene contenido & Archivo faltante o vacio & Critico \\
\bottomrule
\end{tabular}
\end{table}

\newpage

% ============================================================
% 8. PIPELINE COMPLETO
% ============================================================

\end{tcolorbox}


\begin{table}[H]
\centering
\begin{tabular}{llc}
\toprule
\textbf{Verificacion} & \textbf{Metrica} & \textbf{Umbral} \\
\midrule
Archivo .i generado $>$ 10 KB & Tamano $\leq$ 10 KB & Critico \\
Todas las keywords MCNP presentes & Keyword faltante & Critico \\
Coordenadas TMESH dentro del RPP & TMESH fuera de limites & Critico \\
Tarjetas CUT para modos activos & Modo sin CUT & Critico \\
Todas las celdas tienen material $>$ 0 & Celda con material 0 & Critico \\
FILL tiene $nx \times ny \times nz$ elementos & Cuenta incorrecta & Critico \\
Material M definido para cada celda & Material referenciado sin M & Critico \\
Archivo .src existe y tiene contenido & Archivo faltante o vacio & Critico \\
\bottomrule
\end{tabular}
\end{table}

\newpage

% ============================================================
% 8. PIPELINE COMPLETO
% ============================================================
\section{Pipeline Completo --- Integracion del Modulo 2}
\label{sec:pipeline}

\subsection*{Acronimos usados en esta seccion}
\begin{tabular}{ll}
\toprule
\textbf{Acronimo} & \textbf{Significado} \\
\midrule
FS & Final State (estado final del pipeline) \\
QC & Quality Control \\
MIRD & Medical Internal Radiation Dose \\
SIRT & Selective Internal Radiation Therapy \\
TACE & Transarterial Chemoembolization \\
PRRT & Peptide Receptor Radionuclide Therapy \\
.mrb & Medical Reality Bundle \\
.i & Input file (MCNP input) \\
.o & Output file (MCNP output) \\
\bottomrule
\end{tabular}

\subsection{Contexto}

3Dosim es un sistema de dosimetria personalizada para terapias con radionuclidos. El pipeline completo consta de 3 modulos secuenciales:

\begin{verbatim}
  Modulo 1          Modulo 2           Modulo 3
  Segmentacion      Generacion MC      Ejecucion MC
  ----------        ------------       -----------
  TC + PET          Escena .mrb        Input .i
       |                 |                  |
       v                 v                  v
  TotalSegmentator   pipeline_mod2.py    mcnp6
  Labelmap 3D        Input MCNP (.i)    MCTAL
  VOIs .vtk          Espectro (.dat)    Dosis (Gy)
  Escena .mrb        Fuente (.src)      Isodosis
       |                 |                  |
       +-----------------+------------------+
                          |
                     Dosis por VOI
                     Mapas de dosis 3D
                     DVH (histograma dosis-volumen)
\end{verbatim}

\subsection{Interfaces}

\subsubsection{Entrada (desde Modulo 1)}

\begin{center}
\begin{tabular}{lll}
\toprule
\textbf{Elemento} & \textbf{Formato} & \textbf{Descripcion} \\
\midrule
Escena 3D Slicer & .mrb & TC + labelmap + VOIs \\
Labelmap segmentado & NRRD & Indices de 1 a 153 \\
VOIs & VTK & Mallas de higado, tumor, rinones \\
Parametros clinicos & JSON/dict & Isotopo, actividad, $n_{particulas}$ \\
\bottomrule
\end{tabular}
\end{center}

\subsubsection{Salida (hacia Modulo 3)}

\begin{center}
\begin{tabular}{lll}
\toprule
\textbf{Elemento} & \textbf{Formato} & \textbf{Descripcion} \\
\midrule
Input MCNP & .i & Archivo de entrada listo para MCNP6 \\
Fuente externa & .src & Posiciones de particulas muestreadas \\
Espectro & .dat & Distribucion de energias \\
Metadata & .json & Parametros usados en la generacion \\
\bottomrule
\end{tabular}
\end{center}

\subsection{Diagrama de Flujo Detallado}

\begin{verbatim}
                            +-------------------+
                            |     Inicio        |
                            | (SlicerDosim)     |
                            +--------+----------+
                                     |
                                     v
                            +-------------------+
                            | Cargar escena     |
                            | loadScene(.mrb)   |
                            +--------+----------+
                                     |
                                     v
                            +-------------------+
                            | scan_nodes()      |
                            | Identificar:      |
                            | - volume_node     |
                            | - labelmap        |
                            | - VOIs            |
                            +--------+----------+
                                     |
                                     v
                            +-------------------+
                            | validate_         |
                            | prereqs()         |
                            | 4 checks          |
                            +--------+----------+
                                     |
                   +-----------------+------------------+
                   v                 v                  v
             iTC OK?          iLabelmap OK?       iVOIs OK?
                   |                 |                  |
                   +--------+--------+---------+--------+
                            |                  |
                            v                  v
                     Continuar            Detener
                            |
                            v
              +-------------------------------+
              | Configurar TissueConfig       |
              | Cargar tissue_config.json     |
              | (59 materiales totales)        |
              +-------------+-----------------+
                            |
                            v
              +-------------------------------+
              | Construir bloques del .i      |
              |                               |
              | +-> Cabecera (metadata)       |
              | +-> Celdas (lattice + RLE)    |
              | +-> Superficies (RPP)         |
              | +-> Datos (MODE/CUT/IMP)      |
              | +-> Materiales (M1..M6+)      |
              | +-> SDEF (espectro + src)     |
              | +-> TMESH (grilla 3D)         |
              | +-> *F8 (detectores VOI)      |
              | +-> F6 (energia total)        |
              | +-> PRINT/PRDMP               |
              | +-> Archivo .src externo      |
              +-------------+-----------------+
                            |
                            v
              +-------------------------------+
              | Escribir paciente.i           |
              | Escribir paciente.src         |
              | Escribir paciente.dat         |
              +-------------+-----------------+
                            |
                            v
              +-------------------------------+
              | validate_mcnp_input()         |
              | 5 verificaciones              |
              +-------------+-----------------+
                            |
                    +-------+--------+
                    v                v
                Valido          Invalido
                    |                |
                    v                v
        +-------------------+  +----------------+
        | Continuar         |  | Corregir y     |
        | a Modulo 3        |  | regenerar      |
        +-------------------+  +----------------+
\end{verbatim}

\subsection{Dependencias entre Secciones del .i}

\begin{verbatim}
  Superficies (RPP) ----------> Celdas (usan numeros de superficie)
        |                              |
        |                              v
        |                     Materiales (M1..M6+)
        |                              |
        v                              v
  TMESH (coordenadas           Celdas (MAT + RHO)
  deben coincidir con RPP)          |
                                     v
                               SDEF (archivo .src)
                                     |
                                     v
                               *F8 / F6 (celdas de VOI)
\end{verbatim}

\subsection{Manejo de Errores y Estados}

\subsubsection{Estados del Pipeline}

\begin{center}
\begin{tabular}{ll}
\toprule
\textbf{Estado} & \textbf{Significado} \\
\midrule
\texttt{INIT} & Pipeline creado, sin ejecutar \\
\texttt{CHECKING\_SLICER} & Verificando Slicer \\
\texttt{LOADING\_SCENE} & Cargando escena .mrb \\
\texttt{SCANNING\_NODES} & Identificando nodos \\
\texttt{VALIDATING} & Validando prerequisitos \\
\texttt{GENERATING} & Generando input MCNP \\
\texttt{VALIDATING\_OUTPUT} & Validando archivo generado \\
\texttt{COMPLETED} & Pipeline terminado exitosamente \\
\texttt{FAILED} & Pipeline fallo \\
\bottomrule
\end{tabular}
\end{center}

\subsubsection{Errores Comunes}

\begin{center}
\begin{tabular}{lll}
\toprule
\textbf{Error} & \textbf{Causa} & \textbf{Solucion} \\
\midrule
\texttt{SceneLoadError} & .mrb corrupto & Regenerar escena en Mod1 \\
\texttt{NodeNotFound} & Falta volume/labelmap & Verificar segmentacion \\
\texttt{DimensionMismatch} & TC $\neq$ labelmap & Reprocesar Mod1 \\
\texttt{TissueConfigError} & JSON mal formado & Verificar tissue\_config.json \\
\texttt{MCNPValidationError} & Input invalido & Revisar validate\_mcnp\_input.py \\
\texttt{SourceSamplingError} & VOI vacio & Verificar segmentacion tumor \\
\bottomrule
\end{tabular}
\end{center}

\subsection{Integracion con el Resto del Sistema}

\begin{verbatim}
                      +-------------------------+
                      |    3Dosim v4.0          |
                      |    Interfaz Slicer      |
                      +-----------+-------------+
                                  |
                 +----------------+----------------+
                 v                v                v
         +-------------+  +-------------+  +-------------+
         | Modulo 1    |  | Modulo 2    |  | Modulo 3    |
         | Segmentac.  |->| Gen. Input  |->| Ejecucion   |
         | (Python)    |  | (Python)    |  | (MCNP + Py) |
         +-------------+  +-------------+  +-------------+
                                 |                  |
                                 v                  v
                        +-------------+  +-------------+
                        |tissue_config|  |   MCTAL     |
                        |.json        |  |   Parser    |
                        |ICRP-110     |  |   Dosis en  |
                        |59 materiales|  |   Gy        |
                        +-------------+  +-------------+
\end{verbatim}

\subsection{Comandos de Ejecucion}

\subsubsection{Desde 3D Slicer (SlicerDosim)}

\begin{verbatim}
def run_dosimetry_from_scene(mrb_path, params):
    # ===== Modulo 1: Segmentacion =====
    scene = slicer.util.loadScene(mrb_path)
    # ...

    # ===== Modulo 2: Generacion de input MCNP =====
    from pipeline_mod2 import PipelineMod2
    mod2 = PipelineMod2()
    mod2.check_slicer()
    mod2.load_scene(mrb_path)
    mod2.scan_nodes()
    mod2.validate_prereqs()
    input_path = mod2.generate_mcnp(params)
    mod2.validate(input_path)

    # ===== Modulo 3: Ejecucion MCNP =====
    # ...
\end{verbatim}

\subsubsection{Desde Linea de Comandos (Standalone)}

\begin{verbatim}
python pipeline_mod2.py --scene paciente.mrb \
    --isotope Y90 \
    --n_particles 1000000 \
    --output_dir ./output
\end{verbatim}

\subsection{Checkpoints del Pipeline}

El sistema mantiene checkpoints persistentes para reanudar ejecuciones interrumpidas:

\begin{verbatim}
output_dir/
+-- checkpoints/
|   +-- pipeline_checkpoint.json
+-- scenes/
|   +-- 3Dosim_scene.mrb
+-- logs/
|   +-- pipeline_mod2.log
+-- paciente.i
+-- paciente.src
+-- paciente.dat
+-- pipeline_results.json  (historial de ejecuciones)
\end{verbatim}


\begin{tcolorbox}[aiSupervisor]

\begin{tcolorbox}[aiSupervisor]
\subsection*{Control de Calidad (AI Supervisor)}

\begin{table}[H]
\centering
\begin{tabular}{llc}
\toprule
\textbf{Verificacion} & \textbf{Metrica} & \textbf{Umbral} \\
\midrule
Flujo completo ejecutado & Todos los pasos completados & Critico \\
Interfaz Mod1 $\to$ Mod2 & Escena .mrb cargable & Critico \\
Interfaz Mod2 $\to$ Mod3 & Archivo .i valido & Critico \\
Tiempo de generacion & $t_{gen} < 5$ min & Advertencia \\
Conservacion de energia & Ratio 0.95--1.05 & Critico \\
Reproducibilidad & Mismo RAND seed produce mismo resultado & Critico \\
\bottomrule
\end{tabular}
\end{table}

\newpage

% ============================================================
% APENDICE A: CONSTANTES FISICAS
% ============================================================
\section*{Apendice: Constantes Fisicas}
\addcontentsline{toc}{section}{Apendice: Constantes Fisicas}

\end{tcolorbox}


\begin{table}[H]
\centering
\begin{tabular}{llc}
\toprule
\textbf{Verificacion} & \textbf{Metrica} & \textbf{Umbral} \\
\midrule
Flujo completo ejecutado & Todos los pasos completados & Critico \\
Interfaz Mod1 $\to$ Mod2 & Escena .mrb cargable & Critico \\
Interfaz Mod2 $\to$ Mod3 & Archivo .i valido & Critico \\
Tiempo de generacion & $t_{gen} < 5$ min & Advertencia \\
Conservacion de energia & Ratio 0.95--1.05 & Critico \\
Reproducibilidad & Mismo RAND seed produce mismo resultado & Critico \\
\bottomrule
\end{tabular}
\end{table}

\newpage

% ============================================================
% APENDICE A: CONSTANTES FISICAS
% ============================================================
\section*{Apendice: Constantes Fisicas}
\addcontentsline{toc}{section}{Apendice: Constantes Fisicas}

\end{tcolorbox}


\begin{table}[H]
\centering
\begin{tabular}{llc}
\toprule
\textbf{Verificacion} & \textbf{Metrica} & \textbf{Umbral} \\
\midrule
Flujo completo ejecutado & Todos los pasos completados & Critico \\
Interfaz Mod1 $\to$ Mod2 & Escena .mrb cargable & Critico \\
Interfaz Mod2 $\to$ Mod3 & Archivo .i valido & Critico \\
Tiempo de generacion & $t_{gen} < 5$ min & Advertencia \\
Conservacion de energia & Ratio 0.95--1.05 & Critico \\
Reproducibilidad & Mismo RAND seed produce mismo resultado & Critico \\
\bottomrule
\end{tabular}
\end{table}

\newpage

% ============================================================
% APENDICE A: CONSTANTES FISICAS
% ============================================================
\section*{Apendice: Constantes Fisicas}
\addcontentsline{toc}{section}{Apendice: Constantes Fisicas}

\end{tcolorbox}


\begin{table}[H]
\centering
\begin{tabular}{llc}
\toprule
\textbf{Verificacion} & \textbf{Metrica} & \textbf{Umbral} \\
\midrule
Flujo completo ejecutado & Todos los pasos completados & Critico \\
Interfaz Mod1 $\to$ Mod2 & Escena .mrb cargable & Critico \\
Interfaz Mod2 $\to$ Mod3 & Archivo .i valido & Critico \\
Tiempo de generacion & $t_{gen} < 5$ min & Advertencia \\
Conservacion de energia & Ratio 0.95--1.05 & Critico \\
Reproducibilidad & Mismo RAND seed produce mismo resultado & Critico \\
\bottomrule
\end{tabular}
\end{table}

\newpage

% ============================================================
% APENDICE A: CONSTANTES FISICAS
% ============================================================
\section*{Apendice: Constantes Fisicas}
\addcontentsline{toc}{section}{Apendice: Constantes Fisicas}

\subsection*{Acronimos usados en esta seccion}
\begin{tabular}{ll}
\toprule
\textbf{Acronimo} & \textbf{Significado} \\
\midrule
MeV & Megaelectronvoltio \\
keV & Kiloelectronvoltio \\
Bq & Becquerel \\
Gy & Gray \\
\bottomrule
\end{tabular}

\subsection{Isotopos}

\begin{center}
\begin{tabular}{lllll}
\toprule
\textbf{Isotopo} & $T_{1/2}$ (dias) & $E_{\beta,max}$ (MeV) & $\bar{E}_\beta$ (MeV) & \textbf{Alcance medio} \\
\midrule
$^{90}$Y & 2.67 & 2.28 & 0.935 & 2.5 mm \\
$^{177}$Lu & 6.65 & 0.50 & 0.133 & 0.5 mm \\
$^{131}$I & 8.02 & 0.61 & 0.192 & 0.8 mm \\
\bottomrule
\end{tabular}
\end{center}

\subsection{Densidades de Referencia}

\begin{center}
\begin{tabular}{lll}
\toprule
\textbf{Tejido} & $\rho$ (g/cm$^3$) & \textbf{Fuente} \\
\midrule
Aire seco & 0.001205 & ICRU-44 \\
Tejido blando & 1.06 & ICRU-44 \\
Higado & 1.05 & ICRP-110 \\
Pulmon & 0.50 & ICRP-110 \\
Hueso & 1.60 & ICRP-110 \\
Tumor & 1.05 & ICRP-110 \\
Agua & 1.00 & ICRU-44 \\
\bottomrule
\end{tabular}
\end{center}

\subsection{Espectro Y-90 (Detallado)}

\begin{center}
\begin{tabular}{ll}
\toprule
\textbf{Parametro} & \textbf{Valor} \\
\midrule
$E_{\beta,max}$ & 2.2807 MeV \\
$\bar{E}_\beta$ & 0.9357 MeV \\
Decaimiento & $\beta^-$ puro a $^{90}$Zr (estable) \\
$Z_{hijo}$ & 40 ($^{90}$Zr) \\
$T_{1/2}$ & 64.04 horas (2.67 dias) \\
$\lambda$ & $1.082 \times 10^{-2}$ h$^{-1}$ \\
$A_{tipica}$ SIRT & 1--5 GBq \\
\bottomrule
\end{tabular}
\end{center}

\subsection{Factores de Conversion}

\begin{center}
\begin{tabular}{lll}
\toprule
\textbf{De} & \textbf{A} & \textbf{Factor} \\
\midrule
MeV/g & Gy & $1.602176634 \times 10^{-10}$ \\
mCi & MBq & 37 \\
Ci & GBq & 37 \\
Bq & Ci & $2.7027 \times 10^{-11}$ \\
shakes & s & $10^{-8}$ \\
MeV & J & $1.602176634 \times 10^{-13}$ \\
\bottomrule
\end{tabular}
\end{center}

\newpage

% ============================================================
% REFERENCIAS
% ============================================================
\section*{Referencias}
\addcontentsline{toc}{section}{Referencias}

\begin{enumerate}
    \item MCNP6 User's Manual, LA-UR-13-22934, Los Alamos National Laboratory, 2013.
    \item MCNP6 LATTICE Card, LA-CP-13-00634, Los Alamos National Laboratory, 2013.
    \item MCNP6 TMESH Card, LA-CP-14-00745, Los Alamos National Laboratory, 2014.
    \item ICRP Publication 110: Adult Reference Computational Phantoms, Annals of the ICRP 39(2), 2009.
    \item ICRU Report 44: Tissue Substitutes in Radiation Dosimetry and Measurement, 1989.
    \item ICRU Report 56: Dosimetry of External Beta Rays for Radiation Protection, 1997.
    \item Wasserthal, J. et al., ``TotalSegmentator: Robust Segmentation of 104 Anatomic Structures in CT Images'', Radiology Insights, 2023.
    \item Fedorov, A. et al., ``3D Slicer as an Image Computing Platform for the Quantitative Imaging Network'', Magnetic Resonance Imaging 30(9), 2012.
    \item X-5 Monte Carlo Team, ``MCNP -- A General N-Particle Transport Code, Version 5'', LA-UR-03-1987, 2003.
    \item Eckerman, K. et al., ``Nuclear Decay Data for Dosimetric Calculations'', ICRP Publication 107, Annals of the ICRP 38(3), 2008.
    \item Loevinger, R. et al., ``MIRD Primer for Absorbed Dose Calculations'', Society of Nuclear Medicine, 1988.
    \item Schneider, W. et al., ``Correlation between CT numbers and tissue parameters needed for Monte Carlo simulations'', Physics in Medicine and Biology 45(2), 2000.
    \item Solberg, T. et al., ``Lattice and voxel geometry in MCNP for medical physics applications'', Nuclear Science and Engineering, 2001.
    \item Shultis, J.K. \& Faw, R.E., ``An MCNP Primer'', Kansas State University, 2011.
    \item 3Dosim v4 Architecture, \texttt{docs/modulo\_1/vision\_general.md}, 2026.
    \item 3Dosim v4 PipelineOrchestrator, \texttt{pipeline\_mod2.py}, 2026.
    \item 3Dosim v4 SlicerDosim, \texttt{SlicerDosimMod3.py}, 2026.
    \item 3D Slicer MRML Documentation, \texttt{mrml.vtk.org}.
    \item .mrb File Format Specification, 3D Slicer Documentation.
\end{enumerate}

\end{document}
